\chapter{Assembly and consequences of the envelope theorem}
\index{envelope theorem!consequences}\index{unibranch variety!characterization}
\label{chap:bj-envelope-consequences}

All descent mechanisms are now available.
The normal case was proved by alterations in
\Cref{chap:bj-alterations}.
The normalization of a unibranch variety is a finite birational universal
homeomorphism by \Cref{lem:bj-unibranch-normalization-universal}.
Pullback through such a map identifies psh spaces in every characteristic by
\Cref{cor:bj-universal-reflection-all-characteristics}.
This chapter assembles those statements, proves the theorem announced in
\Cref{chap:bj-envelope-statement}, and records the formal consequences
that will be used in the energy chapters.

Throughout, $k$ is algebraically closed and trivially valued,
$X$ is an integral projective $k$-variety, and numerical classes are real.
Extended-real values are retained on the whole of $X^{\An}$.

\section{Orientation}
\label{sec:ch45-orientation}

\subsection{What is proved here, and out of what}
\label{sec:ch45-what}

\Crefrange{chap:bj-cartier-realization}{chap:bj-subintegral-normalization}
proved eight independent theorems, each of which is a statement about
\emph{one} class of morphisms and \emph{one} technique.  None of them is the
envelope theorem.  This chapter is where they are put together, and where
the result is allowed to act on the rest of the book.

The assembly itself occupies four lines and is written out in
\Cref{sec:ch45-proof}.  It reproduces the argument of \S7 of the supplied
manuscript, in which the chain
\[
 \text{smooth}
 \ \xRightarrow{\ \text{alterations}\ }\ \text{normal}
 \ \xRightarrow{\ \text{normalization}\ }\ \text{unibranch}
\]
is closed by a single application of the identification of analytic spaces
under a finite birational universal homeomorphism.  What deserves attention
is not the length of the argument but three things that are easy to get
wrong and that the chapter checks explicitly:

\begin{itemize}
\item that the assembly consumes its inputs \emph{in one direction only},
      so that no input was itself proved using the conclusion
      (\Cref{sec:ch45-noncircular});
\item that the identity of regularized suprema transported through the
      normalization is an identity of \emph{extended-real} functions, valid
      at semivaluations of proper support where both sides may be
      $-\infty$ (\Cref{sec:ch45-endpoint});
\item that the class $\theta$ is an arbitrary element of
      $N^1(X)_{\mathbf R}$, so that the reduction to rational or positive
      classes has to be performed, not assumed
      (\Cref{sec:ch45-perturbation}).
\end{itemize}

The consequences occupy the bulk of the chapter and are of four kinds.
First, the increasing-net formulation
(\Cref{sec:ch45-nets}), which is the form in which the envelope property is
used in every later argument, because psh functions are defined by
\emph{decreasing} approximation and every construction that produces a
supremum produces it as an increasing net.  Second, the obstacle envelope
$P_\theta$ (\Cref{sec:ch45-obstacles}), which becomes an operator with
values in $\PSH_X(\theta)$ only at this point in the book; before it, the
symbol $P_\theta$ denoted a function of unknown regularity, and
\Cref{chap:bj-envelope-history} carried an explicit warning to that effect.
Third, the monotone approximation statements
(\Cref{sec:ch45-monotone}) that convert the closure property into a
calculus.  Fourth, the geometric characterization
(\Cref{sec:ch45-characterization}): in an ample class the envelope property
holds \emph{if and only if} $X$ is unibranch, so the hypothesis of the
theorem is not an artefact of the proof.

\subsection{The firewall inverts in this chapter}
\label{sec:ch45-firewall-inverts}

Every chapter from \Cref{chap:bj-cartier-realization} to
\Cref{chap:bj-subintegral-normalization} carried a warning forbidding the
use of \Cref{thm:bj-final-envelope}; \Cref{warn:ch42-firewall} is the
sharpest of them, because on a normal variety the forbidden appeal would
have shortened a whole chapter to a line.  Those warnings are discharged
here.

\begin{important}{The firewall is discharged in this chapter}
\label{imp:ch45-firewall-discharged}
From this point onward in the book, \Cref{thm:bj-final-envelope} is
available and may be used freely on integral projective unibranch varieties
over an algebraically closed trivially valued field.  In particular:
\begin{enumerate}
\item the envelope property $\operatorname{EP}_X(\theta)$ of
      \Cref{def:bj-envelope-property} may be quoted as a theorem, for every
      $\theta\in N^1(X)_{\mathbf R}$, whenever $X$ is unibranch;
\item $\PSH_X(\theta)$ may be used as a class closed under regularized
      suprema of nonempty bounded-above families;
\item the obstacle envelope $P_\theta$ of
      \Cref{def:bj-obstacle-envelope} may be treated as an operator with
      values in $\PSH_X(\theta)$, subject to the hypotheses recorded in
      \Cref{thm:bj-final-obstacle-psh};
\item the equivalences of \Cref{thm:bj-envelope-equivalences} and
      \Cref{thm:general-obstacle-criterion}, and the increasing-net
      criterion \Cref{prop:bj-ep-equivalences}, may be used in either
      direction without further hypothesis.
\end{enumerate}
The restriction to unibranch $X$ is not a residue of the method:
by \Cref{thm:bj-final-unibranch-characterization} it is equivalent to the
envelope property in a single ample class.
\end{important}

The inversion is worth stating in the negative as well.  Nothing proved
before this chapter may be re-proved using
\Cref{thm:bj-final-envelope}; the earlier proofs are the reason the theorem
is available, and replacing one of them by an appeal to the theorem would
turn the chain into a circle.  \Cref{sec:ch45-noncircular} carries out that
audit once and for all, in the form in which a reader can check it: for each
input, the chapter in which it is proved and the fact that its proof uses
only strictly earlier material.

\subsection{Standing hypotheses and notation}
\label{sec:ch45-hypotheses}

Unless a statement says otherwise:

\begin{itemize}
\item $k$ is algebraically closed and carries the trivial absolute value;
\item $X$ is an integral projective $k$-variety;
\item $\theta\in N^1(X)_{\mathbf R}$ is an arbitrary real numerical class;
\item $\PSH_X(\theta)$ is the Boucksom--Jonsson class of
      \Cref{def:bj-psh}, whose members are upper semicontinuous, generically
      finite, and decreasing limits of nets of generically finite
      Fubini--Study functions in classes \emph{converging} to $\theta$;
\item $X^{\divis}\subset X^{\An}$ is the set of the trivial valuation
      together with the positive rational multiples of divisorial valuations,
      all of which have full support;
\item $\nu:X^\nu\to X$ is the normalization, and $\nu^{\An}$ its
      analytification.
\end{itemize}

Two conventions from the front matter are in force throughout and are worth
repeating, since the whole chapter is about suprema of extended-real
functions.  For a family $(f_i)$ of functions
$X^{\An}\to\Rbar$, the symbol $(\sup_i f_i)^*$ denotes the least upper
semicontinuous majorant of the pointwise supremum; and no difference of two
infinite values is ever formed.  Every identity below involving $-\infty$
is either an identity of suprema or infima, or an identity obtained by
transporting neighbourhood filters through a homeomorphism; neither
operation requires subtraction.

The word \emph{unibranch} is used in the sense of
\Cref{def:ch43-unibranch}: the normalization $\nu$ is bijective on points.
Over an algebraically closed field this is equivalent, for integral $X$,
to $\nu$ being a universal homeomorphism
(\Cref{lem:bj-unibranch-normalization-universal}).

\subsection{What is used from earlier chapters}
\label{sec:ch45-inputs}

The chapter uses six groups of results proved earlier in the book for the
assembly, and one published converse for the characterization.  We keep
these roles separate in the audit below.

\begin{enumerate}
\item \Cref{thm:bj-smooth-envelope}: the envelope property in every real
      class on a smooth projective variety, in every characteristic
      (\Cref{chap:bj-smooth-characteristic}).
\item \Cref{thm:bj-birational-ep-descent}: descent of the envelope property
      along a projective birational morphism of normal projective varieties
      (\Cref{chap:bj-reflection}).
\item \Cref{thm:bj-finite-normal-descent}: descent along a finite surjective
      morphism of normal projective varieties
      (\Cref{chap:bj-finite-descent}).
\item \Cref{thm:bj-normal-envelope} and
      \Cref{cor:bj-normalization-source-ep}: the envelope property on a
      normal projective variety, and on the normalization of an arbitrary
      integral projective variety (\Cref{chap:bj-alterations}).
\item \Cref{lem:bj-unibranch-normalization-universal} and
      \Cref{cor:bj-universal-usc}: unibranchness makes $\nu$ a finite
      birational universal homeomorphism, and $\nu^{\An}$ a homeomorphism
      compatible with upper semicontinuous regularization
      (\Cref{chap:bj-universal-homeomorphisms}).
\item \Cref{cor:bj-universal-reflection-all-characteristics}: pullback is a
      bijection of psh classes through such a map, in every characteristic
      (\Crefrange{chap:bj-universal-homeomorphisms}{chap:bj-subintegral-normalization}).
\item \Cref{thm:ep-implies-unibranch}: the converse in the ample case, due
      to Boucksom and Jonsson and \emph{not} part of the supplied envelope
      manuscript (\Cref{chap:bj-envelope-history}).
\end{enumerate}

Item~(7) is the only one of the seven groups that is not proved in this book, and
the distinction is maintained throughout: the supplied manuscript proves
sufficiency of unibranchness and says nothing about necessity.

\section{The dependency chain}
\label{sec:ch45-dependency}

The proof is short because its ingredients were proved in the correct order.
Their logical dependencies are
\[
\begin{gathered}
 \text{smooth EP}\Longrightarrow\text{normal EP}
 \Longrightarrow\text{EP on }X^\nu
 \Longrightarrow\text{EP on unibranch }X,\\
 \text{Cartier realization}\Longrightarrow
 \text{birational reflection},\\
 \text{Frobenius or Roberts--Singh descent}\Longrightarrow
 \text{normalization reflection}.
\end{gathered}
\tag{45.1}\label{eq:bj-final-dependency-diagram}
\]
Finite normal descent enters the first row between smooth and normal EP,
through the Stein factorization of an alteration.
The characteristic-$p$ and characteristic-zero normalization arguments
enter the lower-right vertex.

\subsection{The chain, with attributions}
\label{sec:ch45-attributions}

It is worth expanding \eqref{eq:bj-final-dependency-diagram} into the exact
sequence of implications used in \Cref{sec:ch45-proof}, with the chapter of
proof attached to each arrow.  Read from the top, the chain is a sequence of
statements about varieties of decreasing regularity: smooth, then normal,
then unibranch.

\begin{enumerate}
\item \emph{Smooth starting point.}
      \Cref{thm:bj-smooth-envelope} establishes
      $\operatorname{EP}_Y(\eta)$ for smooth projective $Y$ and every real
      $\eta$, using multiplier ideals in characteristic zero or asymptotic
      test ideals in positive characteristic, with a uniform
      $O(m^{-1})$ class error (\Cref{chap:bj-smooth-characteristic}).
\item \emph{Two normal-target descents.}
      For a projective birational map $\beta:Y\to T$ of normal varieties,
      \Cref{thm:bj-birational-ep-descent} descends the property from
      $\beta^*\eta$ to $\eta$ by Cartier realization, homogeneous
      reflection and carriers
      (\Cref{chap:bj-cartier-realization,chap:bj-reflection}).
      For a finite surjective map $h:T\to X$ of normal varieties,
      \Cref{thm:bj-finite-normal-descent} descends it from $h^*\eta$ to
      $\eta$ by analytic openness, flattening and determinant-norm fibre
      averages (\Cref{chap:bj-finite-descent}).
\item \emph{The normal case.}
      De Jong alteration and Stein factorization put the smooth starting
      point and both descents together in \Cref{thm:bj-normal-envelope}.
      Applying it to the normal integral projective variety $X^\nu$ gives
      $\operatorname{EP}_{X^\nu}(\nu^*\theta)$ for every integral
      projective $X$, without assuming unibranchness
      (\Cref{cor:bj-normalization-source-ep}).
\item \emph{Topological transport.}
      If $X$ is unibranch, its finite birational normalization is radicial
      and hence a universal homeomorphism
      (\Cref{lem:bj-unibranch-normalization-universal}).
      The resulting analytic homeomorphism commutes with upper
      semicontinuous regularization (\Cref{cor:bj-universal-usc}).
\item \emph{Algebraic reflection.}
      \Cref{cor:bj-universal-reflection-all-characteristics} identifies
      $\PSH_X(\theta)$ with $\PSH_{X^\nu}(\nu^*\theta)$: the descent
      uses a uniform Frobenius power in characteristic $p$ and the
      Roberts--Singh additive-power sandwich in characteristic zero
      (\Cref{chap:bj-universal-homeomorphisms,chap:bj-subintegral-normalization}).
\end{enumerate}

Two features of this sequence are structural rather than incidental.

First, the regularity of the variety decreases along the chain while the
generality of the class never changes.  Each result just cited
holds for an arbitrary real class in $N^1$ of the variety at hand.  The
reduction to rational classes in the proof of the smooth case is internal to
\Cref{chap:bj-smooth-characteristic} and is discharged there
(\Cref{lem:ch40-ep-nef-limit,lem:ch40-rational-approximation}); it is
recorded again, in the form in which the present chapter uses it, in
\Cref{sec:ch45-perturbation}.

Second, the sequence has two different kinds of input.  The smooth case and
the two descent theorems establish the envelope \emph{property}, which then
holds on the normalization.  The homeomorphism identifies upper
semicontinuous regularizations, while reflection transports psh
\emph{functions}.  The distinction matters, because
the envelope property is a statement about all families at once, while the
reflection statement is a statement about individual functions, and the
assembly needs both: it needs the property upstairs to know that a
particular regularized supremum is psh on $X^\nu$, and it needs the
reflection to bring that one function down to $X$.

\subsection{Why the order cannot be permuted}
\label{sec:ch45-order}

Each of the three implications
\[
 \text{smooth}\Longrightarrow\text{normal}\Longrightarrow\text{unibranch}
\]
uses a technique that is unavailable at the next stage, and this is why they
occur in this order and not another.

\emph{Smooth to normal.}  A normal variety need not admit a resolution in
positive characteristic, so the passage cannot be made by a birational
morphism from a smooth model.  It is made by an \emph{alteration}, which is
generically finite of degree $\delta\ge1$, possibly greater than one, and
is split by Stein factorization into a birational part and a finite part.
Both parts are descents, and both
require normality of source and target: birational descent uses the
b-divisor calculus of \Cref{chap:bj-homogeneity-bdivisors}, which is set up
on normal varieties, and finite descent uses fibre averaging, whose
normalization by the generic degree is meaningful only when the varieties
are normal, hence the pure transform is flat after a Raynaud--Gruson blowup.

\emph{Normal to unibranch.}  The normalization $\nu:X^\nu\to X$ is finite
and birational, so it is a degenerate case of both previous steps, and
neither applies: birational descent needs a normal target, and finite
descent needs a normal target.  $X$ is neither.  What is available instead
is that $\nu$ is a \emph{universal homeomorphism}, and hence that psh
functions can be transported by hand, using the ring-theoretic structure of
the extension $\cO_X\subseteq\nu_*\cO_{X^\nu}$: Frobenius in characteristic
$p$, subintegrality and additive powers in characteristic zero.

\emph{Why not directly.}  One might hope to descend from a smooth alteration
of $X$ itself, bypassing $X^\nu$.  The composite $Y\to X$ with $Y$
smooth is generically finite and surjective, but its Stein factorization
$Y\to T\to X$ has $T$ normal and $T\to X$ finite \emph{surjective with
nonnormal target}, which is exactly the case that
\Cref{chap:bj-finite-descent} does not treat.  Its normal-target hypotheses
cannot simply be dropped from the averaging and reflection argument.  The
factorization through $X^\nu$ separates the proven normal-target descent
from the universal-homeomorphic reflection needed at the nonnormal target.

\subsection{The non-circularity audit}
\label{sec:ch45-noncircular}

No arrow in \eqref{eq:bj-final-dependency-diagram} was proved using
\Cref{thm:bj-final-envelope}.
In particular:
\begin{enumerate}
\item Cartier realization was constructed from ordered semiample
      approximations and PSH2;
\item birational reflection was proved from realization and homogeneous
      carriers before source EP was used to descend EP;
\item the smooth theorem was proved independently by multiplier or test
      ideals;
\item finite descent used flattening, determinant norms, and the independent
      carrier-reflection theorem;
\item universal-homeomorphic reflection used Frobenius or Roberts--Singh
      additive powers, not an envelope.
\end{enumerate}
This is the circularity audit for the final three-line argument.

The audit deserves to be carried out in the form in which it can be checked,
because the shape of the chain makes one particular circle very short: a
normal variety is unibranch, so \Cref{thm:bj-final-envelope} applied to a
normal variety would prove \Cref{thm:bj-normal-envelope} in a line.  The
following proposition records that this never happens.

\begin{proposition}
\label{prop:ch45-noncircular}
Order the chapters of Part~IV by their numbers.  Then each of the first six
groups of inputs in \Cref{sec:ch45-inputs} is established in an earlier
chapter without using a result proved later in Part~IV.  The seventh group
is the separately cited published converse and is not used in the assembly.
In particular no proof of an assembly input uses
\Cref{thm:bj-final-envelope}, \Cref{def:bj-envelope-property-final},
\Cref{cor:bj-historical-conjecture}, or any statement of
\Cref{chap:bj-envelope-consequences}.  The deduction of
\Cref{thm:bj-final-envelope} in \Cref{sec:ch45-proof} uses the first six
groups in the order
\[
 \text{(1), (2), (3)}\to\text{(4)}\to\text{(5), (6)} .
\]
\end{proposition}

\begin{proof}
The verification is a finite check, and we run it item by item using the
firewall warnings that each chapter carries.

\emph{(1) \Cref{thm:bj-smooth-envelope}.}  Its proof
(\Cref{chap:bj-smooth-characteristic}) begins with the reduction of
\Cref{lem:ch40-ep-nef-limit} and \Cref{lem:ch40-rational-approximation},
which use only \Cref{thm:bj-basic-closure} and
\Cref{lem:ch34-psh-nonempty} from \Cref{chap:bj-fs-pl-psh}; it then proves
that $P_L(\phi_D)$ is psh for a rational PL obstacle on a big rational
class, by base ideals and multiplier or asymptotic test ideals.  The
envelope property occurs in that chapter as a \emph{conclusion} only.  The
obstacle criterion it uses, \Cref{thm:general-obstacle-criterion}, is proved
in \Cref{chap:bj-envelope-history} from the definition of $\PSH$ and the
comparison interface, with no envelope input.

\emph{(2) \Cref{thm:bj-birational-ep-descent}.}  Its proof
(\Cref{chap:bj-reflection}) uses \Cref{thm:bj-real-cartier-realization} from
\Cref{chap:bj-cartier-realization}, the homogeneous reflection theorem
\Cref{thm:bj-homogeneous-reflection}, the carrier theorem
\Cref{thm:bj-carrier-reflection}, and the canonical carrier
\Cref{lem:bj-canonical-carrier}.  Of these only the last mentions an
envelope, and it mentions it as a \emph{hypothesis on the source}: the
statement is that if the source has the envelope property, then the envelope
of all sup-normalized psh functions is a homogeneous carrier.  It is never
applied to a variety whose envelope property has not been established
independently by an earlier item of the present list.

\emph{(3) \Cref{thm:bj-finite-normal-descent}.}  Its proof
(\Cref{chap:bj-finite-descent}) uses analytic openness
(\Cref{lem:bj-finite-normal-open}), flattening
(\Cref{lem:bj-finite-flatification}), the coefficient-norm construction
(\Cref{prop:bj-weighted-coefficient-norm}), and the carrier theorem of
\Cref{chap:bj-reflection}.  As in the previous item, the envelope property
enters only as a hypothesis on the source of the map.

\emph{(4) \Cref{thm:bj-normal-envelope} and
\Cref{cor:bj-normalization-source-ep}.}  The normal-case proof
(\Cref{chap:bj-alterations}) uses (1), (2) and (3) and nothing else of
substance; \Cref{warn:ch42-firewall} lists, item by item, what may not be
used there, and \Cref{prop:bj-alteration-hypothesis-audit} verifies that
none of it is.  The source of the birational descent is the smooth variety
$Y$, whose envelope property comes from (1); the source of the finite
descent is the normal variety $T$, whose envelope property comes from (2)
applied to $Y\to T$.  Neither appeal is to the theorem being proved.
The normalization-source corollary applies this normal-case theorem to
$X^\nu$, which is normal, integral and projective.  It does not assume
that $X$ is unibranch.

\emph{(5) \Cref{lem:bj-unibranch-normalization-universal} and
\Cref{cor:bj-universal-usc}.}  The first proof uses unibranchness,
finiteness and the diagonal criterion to make the normalization a universal
homeomorphism.  The second is the neighbourhood-filter identity for upper
semicontinuous regularization under the resulting analytic homeomorphism;
\Cref{cor:ch43-sup-star} combines it with pointwise suprema.  These are
algebraic and topological facts proved in
\Cref{chap:bj-universal-homeomorphisms}, without forming a psh envelope.

\emph{(6) \Cref{cor:bj-universal-reflection-all-characteristics}.}  Its
positive-characteristic proof uses a single uniform Frobenius power
(\Cref{lem:bj-uniform-frobenius},
\Cref{cor:ch43-frobenius-all-affines}) carrying every local section upstairs
into $\cO_X$.  In characteristic zero it uses the two-sided valuative sandwich
of \Cref{prop:ch44-global-sandwich}, obtained from the additive-power
theorem \Cref{thm:appA-roberts-singh}.  Neither argument forms an envelope,
and \Cref{prop:ch44-no-birational} records that the characteristic-zero
argument does not even use birationality.

\emph{(7) \Cref{thm:ep-implies-unibranch}.}  This is quoted from
\cite[Theorem~5.12]{BJ22} and is not used in the proof of
\Cref{thm:bj-final-envelope}; it is used only in
\Cref{thm:bj-final-unibranch-characterization}, and there only in the
direction opposite to the theorem.

Finally, the deduction in \Cref{sec:ch45-proof} uses (4) to obtain the
envelope property on $X^\nu$, (5) to identify the two analytic spaces and
their regularized suprema, and (6) to reflect the resulting psh function.
Item (7) enters only in the later characterization.  Thus the assembly
uses exactly the inputs named in \Cref{sec:ch45-inputs}, without a backward
edge in the dependency chain.
\end{proof}

\begin{remark}[What acyclicity does and does not certify]
\label{rem:ch45-audit-scope}
\Cref{prop:ch45-noncircular} certifies that the \emph{book's} proof of
\Cref{thm:bj-final-envelope} is not circular.  It does not certify that the
individual imported statements are correct, and it does not certify that the
external inputs quoted inside them --- de Jong's alterations, Raynaud--Gruson
flattening, Canton's log-discrepancy bound, Musta\c t\u a's uniform
generation, the Roberts--Singh additive-power theorem --- are independent of
non-Archimedean pluripotential theory.  They are, evidently; but the
verification of that is a matter of reading their sources, not of the
ordering used here.  The audit is a statement about the internal edges of the
dependency graph, and it is exactly what a reader is entitled to demand of a
nine-chapter proof whose final step is four lines long.
\end{remark}

\section{Proof of the envelope theorem}
\label{sec:ch45-proof}

\Cref{chap:bj-envelope-statement} already gave the argument, in the
compressed form appropriate to a chapter whose purpose was to state the
theorem.  We repeat it here at the level of detail appropriate to a chapter
whose purpose is to close the proof: every hypothesis of every input is
checked, and every identity of extended-real functions is justified by the
operation that produces it.

\subsection{The three inputs, with their hypotheses verified}
\label{sec:ch45-three-inputs}

Fix an integral projective unibranch $k$-variety $X$ and a class
$\theta\in N^1(X)_{\mathbf R}$, and let
\[
                         \nu:X^\nu\longrightarrow X
\tag{45.2}\label{eq:bj-final-normalization}
\]
be the normalization.  Three facts are used, and each carries a hypothesis
that must be verified for the data at hand.

\begin{lemma}
\label{lem:ch45-normalization-data}
Let $X$ be an integral projective unibranch variety over an algebraically
closed field $k$, with normalization \eqref{eq:bj-final-normalization}.
Then:
\begin{enumerate}
\item\label{itm:ch45-nd:normal}
      $X^\nu$ is a normal integral projective $k$-variety and $\nu$ is
      finite, birational and surjective;
\item\label{itm:ch45-nd:uh}
      $\nu$ is a universal homeomorphism, and
      $\nu^{\An}:(X^\nu)^{\An}\to X^{\An}$ is a homeomorphism of compact
      Hausdorff spaces;
\item\label{itm:ch45-nd:ep}
      $\operatorname{EP}_{X^\nu}(\nu^*\theta)$ holds for every
      $\theta\in N^1(X)_{\mathbf R}$;
\item\label{itm:ch45-nd:refl}
      $\nu^*:\PSH_X(\theta)\to\PSH_{X^\nu}(\nu^*\theta)$ is a bijection,
      in every characteristic.
\end{enumerate}
\end{lemma}

\begin{proof}
\cref{itm:ch45-nd:normal}.  A scheme of finite type over a field is Nagata,
so its normalization is finite
\cite[\href{https://stacks.math.columbia.edu/tag/035S}{Tag 035S}]
{stacks-project}; a finite morphism is projective, so $X^\nu$ is projective
over $k$.  Normality and integrality hold by construction, birationality
because the normalization does not change the function field, and
surjectivity because a finite birational morphism onto an integral scheme is
surjective.  This is the content of the first paragraph of the proof of
\Cref{cor:bj-normalization-source-ep}.

\cref{itm:ch45-nd:uh}.  This is
\Cref{lem:bj-unibranch-normalization-universal}, whose proof uses exactly the
unibranch hypothesis: unibranchness of the integral $X$ says that $\nu$
is bijective, and over an algebraically closed field a finite surjective
bijective morphism is radicial, hence a universal homeomorphism.  The
analytic consequence is
\Cref{lem:bj-universal-analytic-homeomorphism}: $\nu^{\An}$ is a
homeomorphism.  Compactness and the Hausdorff property of both spaces are
the standing facts of \Cref{chap:trivial-analytification} for projective
varieties.

\cref{itm:ch45-nd:ep}.  This is \Cref{cor:bj-normalization-source-ep}, which
is \Cref{thm:bj-normal-envelope} applied to the normal integral projective
variety $X^\nu$ and the real class $\nu^*\theta$.  Two hypotheses are
worth naming.  The variety $X^\nu$ is normal and projective by
\cref{itm:ch45-nd:normal}, which is what \Cref{thm:bj-normal-envelope}
requires; and $\nu^*\theta$ is a genuine element of
$N^1(X^\nu)_{\mathbf R}$, numerical pullback of $\mathbf R$-Cartier
classes being defined for an arbitrary morphism of projective varieties by
\Cref{prop:appF-mixed-cartier}.  No positivity of $\theta$ is used, and
none is available.

\cref{itm:ch45-nd:refl}.  This is
\Cref{cor:bj-universal-reflection-all-characteristics}, applied to the map
$\nu$, which by \cref{itm:ch45-nd:normal} and \cref{itm:ch45-nd:uh} is a
finite birational universal homeomorphism of integral projective varieties
with \emph{normal source}.  Normality of the source is the hypothesis that
has to be checked, and it holds by \cref{itm:ch45-nd:normal}; it is used in
\Cref{chap:bj-subintegral-normalization} to know that
$\cO_X\subseteq\nu_*\cO_{X^\nu}$ is the normalization inclusion and hence,
in characteristic zero, subintegral
(\Cref{cor:ch44-unibranch-normalization-reflection}\,\cref{itm:ch44-unib:subint}).
\end{proof}

\subsection{The assembly}
\label{sec:ch45-assembly}

\begin{proof}[Proof of \Cref{thm:bj-final-envelope}]
Let $\nu:X^\nu\to X$ be the normalization
\eqref{eq:bj-final-normalization}.
By \Cref{cor:bj-normalization-source-ep}, the normal projective variety
$X^\nu$ has the envelope property in the class $\nu^*\theta$.
By unibranchness and
\Cref{lem:bj-unibranch-normalization-universal}, $\nu$ is a finite
birational universal homeomorphism.
Hence $\nu^{\An}$ is a homeomorphism and
\Cref{cor:bj-universal-reflection-all-characteristics} applies.

Let $(\varphi_i)_{i\in I}\subset\PSH_X(\theta)$ be nonempty and bounded
above, and set
\[
                         u:=\left(\sup_{i\in I}\varphi_i\right)^*.
\tag{45.3}\label{eq:bj-final-envelope-u}
\]
Dominant pullback preserves psh functions, so
$\nu^*\varphi_i\in\PSH_{X^\nu}(\nu^*\theta)$.
The family remains nonempty and bounded above.
Because $\nu^{\An}$ is a homeomorphism,
\Cref{cor:bj-universal-usc} gives
\[
 \begin{aligned}
 \nu^*u
 &=\nu^*\left(\sup_i\varphi_i\right)^*\\
 &=\left(\sup_i\nu^*\varphi_i\right)^*.
 \end{aligned}
\tag{45.4}\label{eq:bj-final-pull-envelope}
\]
The right-hand side belongs to
$\PSH_{X^\nu}(\nu^*\theta)$ by the envelope property on the normalization.
Universal-homeomorphic reflection now yields
\[
                         u\in\PSH_X(\theta).
\]
This is exactly the conclusion of \Cref{thm:bj-final-envelope}.
\end{proof}

\subsection{The same four steps, with each step justified}
\label{sec:ch45-four-steps}

The proof just given is short enough that its four steps can be hidden by
the notation.  We separate them.

\emph{Step 1: the family goes up.}  The morphism $\nu$ is surjective, and
pullback along a surjective morphism of integral projective varieties maps
$\PSH_X(\theta)$ into $\PSH_{X^\nu}(\nu^*\theta)$; this is the
functoriality clause of the psh package, recorded in
\Cref{chap:bj-fs-pl-psh} and used throughout
\Crefrange{chap:bj-reflection}{chap:bj-subintegral-normalization}.  It is
important that this is the \emph{easy} direction: it needs no hypothesis
beyond surjectivity, since it is proved by pulling back an approximating net
of generically finite Fubini--Study functions term by term, and a generically
finite Fubini--Study function pulls back to one because $\nu$ is dominant.
The family $(\nu^*\varphi_i)$ is nonempty because $(\varphi_i)$ is, and
bounded above by the same constant $C$, because $\nu^{\An}$ is
surjective.

\emph{Step 2: the envelope is formed upstairs.}  Here
\Cref{cor:bj-normalization-source-ep} is used, and it is the only place in
the proof where a genuinely hard theorem is consumed.  Its output is that
\[
 v:=\left(\sup_i\nu^*\varphi_i\right)^*
 \in\PSH_{X^\nu}(\nu^*\theta).
\]
Note that the conclusion is about a function on $(X^\nu)^{\An}$, and that
nothing has yet been said about $X$.

\emph{Step 3: the two envelopes are identified.}  This is
\eqref{eq:bj-final-pull-envelope}, and it is a purely topological statement.
Since $\nu^{\An}$ is a homeomorphism, it carries the neighbourhood filter
of a point $y\in(X^\nu)^{\An}$ onto the neighbourhood filter of
$\nu^{\An}(y)$, and both upper semicontinuous regularization and pointwise
supremum are defined by those filters and by the pointwise order.  Hence
$\nu^*$ commutes with both operations, which is
\Cref{cor:bj-universal-usc} and, in the form used here,
\Cref{cor:ch43-sup-star}.  No finiteness is required: the identity is an
identity of $\Rbar$-valued functions.  This is the step in which the
unibranch hypothesis is genuinely consumed for the \emph{topology}; the
following step consumes it for the \emph{algebra}.

\emph{Step 4: the envelope comes down.}  By Steps 2 and 3 the function
$\nu^*u$ lies in $\PSH_{X^\nu}(\nu^*\theta)$.  By
\Cref{cor:bj-universal-reflection-all-characteristics} the map $\nu^*$ is a
bijection onto $\PSH_{X^\nu}(\nu^*\theta)$, so there is a unique
$w\in\PSH_X(\theta)$ with $\nu^*w=\nu^*u$.  Since $\nu^{\An}$ is
surjective, $w=u$ pointwise on $X^{\An}$; hence $u\in\PSH_X(\theta)$.

\begin{remark}[Where each hypothesis of the theorem is used]
\label{rem:ch45-hypothesis-use}
It is instructive to see the four hypotheses of
\Cref{thm:bj-final-envelope} being consumed.
\emph{Unibranchness} is used twice, in Steps 3 and 4, and only there.
\emph{Projectivity} is used in Step 2, through
\Cref{thm:bj-normal-envelope}, and again implicitly in the compactness of
$X^{\An}$, which is what makes the bounded-above hypothesis meaningful.
\emph{Nonemptiness} is used in Step 2, since the envelope property is a
statement about nonempty families, and again in Step 4, where a psh function
must not be identically $-\infty$.
\emph{Boundedness above by a constant} is used in Step 2 only.  The
characteristic of $k$ is used nowhere in this chapter: it is used inside
\Cref{cor:bj-universal-reflection-all-characteristics}, where the two
characteristics receive different proofs, and the corollary is stated so
that the case distinction does not resurface here.
\end{remark}

The proof applies to every real numerical class.
No rationality, bigness, nefness, or ampleness assumption on $\theta$ was
introduced.
If the family in \eqref{eq:bj-final-envelope-u} exists, the psh space is
automatically nonempty; if the psh space is empty, the envelope property is
vacuous by definition.

\begin{corollary}\label{cor:bj-conjecture-514}
For every integral projective unibranch variety over an algebraically closed
trivially valued field,
\[
                         \operatorname{EP}_X(\theta)
\]
holds for every $\theta\in N^1(X)_{\mathbf R}$.
Thus Conjecture~5.14 of \cite{BJ22} holds in this setting.
\end{corollary}

\begin{proof}
Apply \Cref{thm:bj-final-envelope} separately to each real numerical class.
The conjecture is precisely the class-by-class envelope assertion.
\end{proof}

\begin{remark}[The empty case is not a gap]
\label{rem:ch45-empty-case}
\Cref{cor:bj-conjecture-514} asserts $\operatorname{EP}_X(\theta)$ for
\emph{every} real class, including classes for which
$\PSH_X(\theta)=\emptyset$.  There is no sleight of hand: by
\Cref{def:bj-envelope-property} the assertion quantifies over nonempty
families, and if the psh class is empty there are none, so the assertion
holds vacuously.  By \Cref{lem:ch34-psh-nonempty} the psh class is empty
whenever $\theta$ is not pseudoeffective, so the complement of
$\operatorname{Psef}(X)$ lies in the vacuous range.  The converse
nonemptiness assertion is not supplied by that lemma; inside
$\operatorname{Psef}(X)$ the envelope assertion has content whenever the
psh class is nonempty, whether or not the class is nef.  This is the point of
\Cref{ex:ch37-empty-class}, and it is the reason
\Cref{thm:bj-final-envelope} carries no positivity hypothesis.
\end{remark}

\section{The proper-support endpoint}
\label{sec:ch45-endpoint}

The equality \eqref{eq:bj-final-pull-envelope} is an equality of
extended-real functions on the whole analytification.
Suppose $x\in X^{\An}$ has support $Z\subsetneq X$, and let
$y\in(X^\nu)^{\An}$ be its unique lift.
If every function in the family has value $-\infty$ at $x$, then
\[
           \sup_i\varphi_i(x)=-\infty
           \quad\text{and}\quad
           \sup_i\nu^*\varphi_i(y)=-\infty.
\]
Upper semicontinuous regularization is compared by the neighborhood filters
under the homeomorphism; no subtraction occurs.
Thus \eqref{eq:bj-final-pull-envelope} remains true whether its common value
is finite or $-\infty$.

The earlier construction stages protect the same endpoint in four different
ways:
\begin{enumerate}
\item homogeneous Siu division gives a global psh identity in birational
      reflection;
\item ideal equalities preserve common zeros in the smooth argument;
\item the zero determinant polynomial preserves common zeros in finite
      descent;
\item Frobenius powers or the two-sided weighted-ideal sandwich preserve the
      common-zero locus in normalization reflection.
\end{enumerate}
The final homeomorphism merely assembles these already exact identities.
It does not repair a function first defined only on divisorial valuations.

\subsection{Why the endpoint is not a technicality}
\label{sec:ch45-endpoint-why}

It is tempting to regard the values at proper-support semivaluations as a
boundary phenomenon that could be handled by fiat --- define the envelope on
$X^{\divis}$, where everything is finite, and extend by upper
semicontinuity.  This does not work, and it is worth saying precisely why,
because the whole extended-real discipline of
\Crefrange{chap:bj-cartier-realization}{chap:bj-subintegral-normalization}
exists to avoid it.

The obstruction is that the extension by upper semicontinuity is not, a
priori, psh.  A $\theta$-psh function is by \Cref{def:bj-psh} a decreasing
limit of a net of generically finite Fubini--Study functions, and such a
function has a definite value at every point of $X^{\An}$, computed from
the vanishing orders of its defining sections by
\eqref{eq:divisor-section-pole}.  Two psh functions with the same divisorial
trace do agree everywhere, by the comparison interface
\cite[Theorem~4.22]{BJ22} recorded in \Cref{chap:bj-envelope-statement}; but
that statement presupposes that both functions are psh.  It says nothing
about a function manufactured from a divisorial trace by regularization, and
in fact it is exactly the assertion that such a manufactured function is psh
that the envelope theorem supplies.  Using the comparison interface to
\emph{define} the extension would therefore assume the conclusion.

The situation is not hypothetical.  \Cref{ex:gf-proper-support-pole}
exhibits a generically finite Fubini--Study function on $\mathbf P^1$ whose
pole locus is the single point $v_{0,\triv}$, and
\Cref{ex:ch34-decreasing-to-pole} exhibits a decreasing sequence of
finite-valued Fubini--Study functions acquiring that pole in the limit.  A
family of such functions, indexed so that the pole locus moves, has a
supremum whose behaviour at proper-support points is genuinely new
information, not determined by the divisorial trace of the individual
members.

\begin{proposition}\label{prop:bj-final-no-hidden-finiteness}
The proof of \Cref{thm:bj-final-envelope} never assumes that a GF
regularization, its limit, or the envelope is finite on all of
$X^{\An}$.
It uses finiteness only at the trivial valuation when PSH2 must exclude its
identically $-\infty$ alternative.
\end{proposition}

\begin{proof}
Every stagewise identity listed above is meaningful in
$\mathbf R\cup\{-\infty\}$.
The trivial valuation has full support, and every nonzero GF coefficient has
finite value there.
A nonempty psh family also supplies a finite lower witness at that point.
These are exactly the nontriviality checks used in the decreasing-limit
arguments.

In detail, the four occurrences of a decreasing limit in the chain are the
following, and each is accompanied by a finite value at $v_{X,\triv}$.
In \Cref{chap:bj-cartier-realization} the ordered semiample approximations
of \Cref{lem:bj-ordered-rational-approximation} are homogeneous, hence
vanish at $v_{X,\triv}$ by the homogeneity package of
\Cref{chap:bj-homogeneity-bdivisors}; the limit therefore has value $0$
there and is not identically $-\infty$.
In \Cref{chap:bj-smooth-characteristic} the base-ideal stages
$\phi_D+m^{-1}\phi_{\mathfrak a_m}$ have value $\phi_D(v_{X,\triv})$ at
the trivial valuation, because the ideals $\mathfrak a_m$ are nonzero and
therefore have a generator not vanishing on the generic fibre.
In \Cref{chap:bj-finite-descent} the descended function is a positively
weighted average of values of the given function, and the given function is
finite at the trivial valuation.
In \Cref{chap:bj-subintegral-normalization} the sandwich
\eqref{eq:bj-char-zero-gf-approximants} is two-sided with a finite width, so
finiteness at $v_{X,\triv}$ upstairs forces it downstairs.

Finally, the assembly of \Cref{sec:ch45-assembly} contains no decreasing
limit at all: it forms one supremum, regularizes once, transports the result
through a homeomorphism, and applies a bijection.  A fortiori it introduces
no finiteness hypothesis.
\end{proof}

\begin{proposition}
\label{prop:ch45-envelope-can-be-minus-infinity}
There exist $X$, $\theta$, and a nonempty bounded-above family
$(\varphi_i)\subset\PSH_X(\theta)$ on a unibranch $X$ such that
$\bigl(\sup_i\varphi_i\bigr)^*$ takes the value $-\infty$ at some point
of $X^{\An}$.
\end{proposition}

\begin{proof}
Take $X=\mathbf P^1$, which is smooth, hence unibranch, and
$\theta=c_1(\cO(1))$.  Let $s\in H^0(X,\cO(1))$ vanish exactly at
$p=0$ and put $\varphi:=\log|s|$, a generically finite Fubini--Study
function, hence an element of $\PSH_X(\theta)$: the constant net
$\phi_i=\varphi$ satisfies \Cref{def:bj-psh}\,(3), and $\varphi$ is
upper semicontinuous and generically finite by
\Cref{prop:ch34-generic-finiteness}.  By
\Cref{ex:gf-proper-support-pole}, $\varphi(v_{0,\triv})=-\infty$ while
$\varphi\le0$ everywhere.  Apply the theorem to the one-member family
$\{\varphi\}$: it is nonempty and bounded above by $0$, its supremum is
$\varphi$, and $\varphi$ is already upper semicontinuous, so
$\bigl(\sup\bigr)^*=\varphi$, which equals $-\infty$ at
$v_{0,\triv}$.
\end{proof}

The example is trivial but the point it makes is not: the conclusion of
\Cref{thm:bj-final-envelope} is a membership statement in a class of
$\Rbar$-valued functions, and no clause of it asserts finiteness.  A
formulation that quietly restricted attention to real-valued psh functions
would be a different, weaker theorem, and would be useless for the
applications, since the psh functions produced by Fubini--Study
approximation are exactly the ones with poles.

\begin{proposition}
\label{prop:ch45-pole-locus}
Let $X$ be integral projective unibranch, $\theta\in N^1(X)_{\mathbf R}$,
and let $(\varphi_i)_{i\in I}\subset\PSH_X(\theta)$ be nonempty and
bounded above, with $u:=(\sup_i\varphi_i)^*$.  Then
\[
 \{u=-\infty\}\ \subseteq\ \bigcap_{i\in I}\{\varphi_i=-\infty\},
\]
and the right-hand side is disjoint from $X^{\lin}$.  If moreover some
$\varphi_{i_0}$ is finite at a point $x$, then $u(x)>-\infty$.
\end{proposition}

\begin{proof}
The pointwise supremum dominates each member, so
$\sup_i\varphi_i\ge\varphi_{i_0}$ for every $i_0$, and upper
semicontinuous regularization is order preserving; hence
$u\ge\varphi_{i_0}$.  If $u(x)=-\infty$ then
$\varphi_{i_0}(x)=-\infty$ for every $i_0$, which is the displayed
inclusion, and the last assertion is its contrapositive.  Each set
$\{\varphi_i=-\infty\}$ is disjoint from $X^{\lin}$, because
\Cref{thm:bj-basic-closure}\,(7), sourced from
\cite[Corollary~4.17]{BJ22}, guarantees finiteness there.  No such
finiteness statement holds in general on all of $X^{\val}$; poles may
also occur at full-support valuations outside $X^{\lin}$.
\end{proof}

\begin{remark}[The intersection may be empty although every member has a pole]
\label{rem:ch45-pole-strict}
On $\mathbf P^1$ with $\theta=c_1(\cO(1))$, take the family
$\{\log|s_q|\}$ indexed by the closed points $q$, where $s_q$ vanishes
exactly at $q$.  Every member has a pole, at $v_{q,\triv}$; but at a
fixed $v_{q_0,\triv}$ every member with $q\ne q_0$ takes the value
$0$, so the intersection in \Cref{prop:ch45-pole-locus} is empty and the
supremum is the constant $0$.  Thus a family of psh functions all of which
have poles can have an envelope with no pole at all.
\end{remark}

The inclusion of \Cref{prop:ch45-pole-locus} is strict in general, and the
mechanism by which it becomes strict is exactly the regularization step.
The next example makes this concrete, and it is the promised
psh counterpart of the elementary \Cref{ex:ch36-sup-not-usc}.

\begin{example}[A psh family whose raw supremum is not upper semicontinuous]
\label{ex:ch45-raw-sup-not-usc}
Let $X=\mathbf P^1$ with $\theta=c_1(\cO(1))$, so $\deg\theta=1$, and
let $p=0$ be a closed point.  For $a\in(0,1]$ let
$\varphi_a:X^{\An}\to\mathbf R\cup\{-\infty\}$ be the function
\[
 \varphi_a(t\ord_p)=-at\ \ (t\in[0,+\infty)),
 \qquad
 \varphi_a\equiv0\ \text{ off the ray through }\ord_p .
\]
By the curve criterion recorded at the end of
\Cref{sec:ch34-pl-positivity}, $\varphi_a\in\PSH_X(\theta)$: the
restriction to each ray is affine, hence convex, and the slope condition
reads
\[
 \deg\theta+\sum_q\varphi_a'{}_{,q}(0^+)=1-a\ \ge\ 0 .
\]
For $a=1$ the function is $\log|s_p|$ of
\Cref{ex:gf-proper-support-pole}; for $a<1$ it is $\theta$-psh by the
criterion, although it need not be a Fubini--Study function for
$\cO(1)$ itself.  Each $\varphi_a$ is bounded above by $0$, and each
satisfies
\[
 \varphi_a(v_{p,\triv})=\lim_{t\to+\infty}(-at)=-\infty .
\]
Now compute the pointwise supremum over the family $(\varphi_a)_{a\in(0,1]}$:
\[
 \sup_{a\in(0,1]}\varphi_a(t\ord_p)
 =\sup_{a\in(0,1]}(-at)=0
 \quad\text{for every }t\in[0,+\infty),
\]
the supremum being approached as $a\downarrow0$ and not attained, while
\[
 \sup_{a\in(0,1]}\varphi_a(v_{p,\triv})=-\infty .
\]
So the raw supremum equals $0$ on all of $X^{\An}$ except at the single
point $v_{p,\triv}$, where it equals $-\infty$.  It is therefore
\emph{not} upper semicontinuous: the points $t\ord_p$ converge to
$v_{p,\triv}$ as $t\to+\infty$, and along them the raw supremum is
constantly $0$.  Its upper semicontinuous regularization is the constant
$0$, which is $\theta$-psh because $\theta$ is nef
(\Cref{lem:ch34-psh-nonempty}).

Three lessons are recorded.  First, the star in
\eqref{eq:bj-envelope-definition} is not cosmetic even for psh families on a
smooth rational curve.  Second, the inclusion of
\Cref{prop:ch45-pole-locus} is strict here: every member is $-\infty$ at
$v_{p,\triv}$, while the envelope is $0$ there.  Third, the failure of
upper semicontinuity occurs precisely at a semivaluation of proper support,
which is why the extended-real discipline cannot be confined to a final
remark.
\end{example}

\section{Increasing nets and regularized suprema}
\label{sec:ch45-nets}

The envelope theorem is equivalent to a monotone statement, but nets rather
than sequences are required.
No countable-cofinality property of the relevant directed sets or Berkovich
spaces is assumed.

The reason the monotone form is the useful one is structural.  A
$\theta$-psh function is, by \Cref{def:bj-psh}, a \emph{decreasing} limit;
every closure property proved before this chapter, namely
\Cref{thm:bj-basic-closure}, concerns decreasing limits; and a decreasing
limit of upper semicontinuous functions is upper semicontinuous for free.
An increasing family is exactly the case in which none of that is available,
and \Cref{ex:ch45-raw-sup-not-usc} shows that the missing upper
semicontinuity is a real phenomenon and not a defect of the definition.  The
envelope theorem is the statement that after the regularization has been
performed, one is back inside the class.

\subsection{The monotone formulation}
\label{sec:ch45-monotone-form}

\begin{theorem}\label{thm:bj-final-increasing-nets}
Let $X$ be integral, projective, and unibranch, and fix
$\theta\in N^1(X)_{\mathbf R}$.
If $(u_j)_{j\in J}\subset\PSH_X(\theta)$ is an upward-directed net with a
common finite upper bound, then there is a unique
$u\in\PSH_X(\theta)$ whose divisorial trace is
\[
                         u(v)=\sup_j u_j(v)
               \qquad(v\in X^{\divis}).
\tag{45.5}\label{eq:bj-final-directed-trace}
\]
Moreover,
\[
                         u=\left(\sup_j u_j\right)^*.
\tag{45.6}\label{eq:bj-final-directed-envelope}
\]
\end{theorem}

\begin{proof}
\Cref{thm:bj-final-envelope} gives $\operatorname{EP}_X(\theta)$.
The increasing-net criterion
\cite[Theorem~5.11]{BJ22}, recorded in
\Cref{prop:bj-ep-equivalences}, gives
\eqref{eq:bj-final-directed-trace} and
\eqref{eq:bj-final-directed-envelope}.
Uniqueness follows from psh determination by divisorial values.

We repeat the deduction with the two comparison steps written out, since it
is used constantly below.  Put $w:=(\sup_j u_j)^*$.  The family
$(u_j)_{j\in J}$ is nonempty --- a directed set is nonempty by convention
--- and bounded above by hypothesis, so
$\operatorname{EP}_X(\theta)$ gives $w\in\PSH_X(\theta)$.

\emph{The divisorial trace of $w$.}  Every $u_j$ is $\theta$-psh and
$w$ is $\theta$-psh, and $u_j\le\sup_ju_j\le w$ everywhere; hence
$\sup_ju_j\le w$ on $X^{\divis}$.  Conversely, $\sup_ju_j$ is a
pointwise supremum of an \emph{upward-directed} family of psh functions, so
by the divisorial nonnegligibility recorded in
\Cref{thm:bj-envelope-equivalences} its restriction to $X^{\divis}$
already determines its upper semicontinuous regularization: precisely,
$w=(\sup_ju_j)^*$ and $\sup_ju_j$ agree on $X^{\divis}$.  This gives
\eqref{eq:bj-final-directed-trace} with $u:=w$, and
\eqref{eq:bj-final-directed-envelope} is then the definition of $w$.

\emph{Uniqueness.}  If $u'\in\PSH_X(\theta)$ also satisfies
\eqref{eq:bj-final-directed-trace}, then $u$ and $u'$ agree on
$X^{\divis}$; by the comparison interface \cite[Theorem~4.22]{BJ22},
recorded in \Cref{chap:bj-envelope-statement}, two $\theta$-psh functions
agreeing on $X^{\divis}$ agree on $X^{\An}$.  Hence $u=u'$.
\end{proof}

\begin{corollary}\label{cor:bj-final-finite-maxima}
Let $X$ be integral projective unibranch.  For a nonempty bounded-above
family
$(\varphi_i)_{i\in I}\subset\PSH_X(\theta)$, let
\[
                         u_F:=\max_{i\in F}\varphi_i
\tag{45.7}\label{eq:bj-final-finite-maxima}
\]
for every finite nonempty $F\subset I$.
Then $(u_F)$, ordered by inclusion of $F$, is an increasing net in
$\PSH_X(\theta)$, and its psh limit is
\[
                         \left(\sup_i\varphi_i\right)^*.
\]
\end{corollary}

\begin{proof}
Finite maxima in one class are psh.
The finite nonempty subsets form a directed set, and their raw pointwise
supremum is $\sup_i\varphi_i$.
Apply \Cref{thm:bj-final-increasing-nets}.

For completeness: that finite maxima are $\theta$-psh is
\Cref{thm:bj-basic-closure}\,(2) together with induction on $\#F$; that
the finite nonempty subsets of $I$ form a directed set under inclusion is
immediate, $F\cup F'$ dominating both $F$ and $F'$; that the net is
increasing is monotonicity of the maximum; and that
$\sup_Fu_F=\sup_i\varphi_i$ holds pointwise, since every $\varphi_i$
occurs as $u_{\{i\}}$ and every $u_F$ is a maximum of members of the
family.  This is \Cref{lem:ch36-directed-reduction}, which was proved in
\Cref{chap:bj-envelope-history} without any envelope input; what is new here
is that the limit exists in $\PSH_X(\theta)$.
\end{proof}

\begin{corollary}\label{cor:bj-final-usc-increasing-limit}
In the setting of \Cref{thm:bj-final-increasing-nets}, if the raw pointwise
supremum $\sup_j u_j$ is already upper semicontinuous, then
\[
                         \sup_j u_j\in\PSH_X(\theta).
\]
\end{corollary}

\begin{proof}
An upper semicontinuous function equals its upper semicontinuous
regularization, so use \eqref{eq:bj-final-directed-envelope}.
\end{proof}

\begin{remark}[What the theorem does not say pointwise]
\label{rem:ch45-directed-pointwise}
Without the upper semicontinuity hypothesis in the last corollary, the raw
supremum need not equal its regularization at every proper-support point.
The canonical equality before regularization is the divisorial trace
\eqref{eq:bj-final-directed-trace}.
The global psh function is \eqref{eq:bj-final-directed-envelope}.
\end{remark}

\Cref{ex:ch45-raw-sup-not-usc} is the case in point: there the two functions
differ at exactly one point of $X^{\An}$, a semivaluation of proper
support, and agree on $X^{\divis}$ as
\eqref{eq:bj-final-directed-trace} predicts.

\subsection{The supremum as a least upper bound}
\label{sec:ch45-lub}

The next statement is what one actually uses: the psh function produced by
the theorem is not merely \emph{a} psh function with the right divisorial
trace, it is the least upper bound of the family inside $\PSH_X(\theta)$
for the pointwise order.

\begin{proposition}
\label{prop:ch45-lub}
Let $X$ be integral projective unibranch and
$\theta\in N^1(X)_{\mathbf R}$.  Let
$(\varphi_i)_{i\in I}\subset\PSH_X(\theta)$ be nonempty and bounded above,
and put $u:=(\sup_i\varphi_i)^*$.  Then:
\begin{enumerate}
\item\label{itm:ch45-lub:upper}
      $u\in\PSH_X(\theta)$ and $\varphi_i\le u$ on $X^{\An}$ for every
      $i$;
\item\label{itm:ch45-lub:least}
      if $w\in\PSH_X(\theta)$ satisfies $\varphi_i\le w$ for every
      $i$, then $u\le w$ on $X^{\An}$;
\item\label{itm:ch45-lub:usc}
      more generally, if $h:X^{\An}\to\Rbar$ is upper semicontinuous and
      $\varphi_i\le h$ for every $i$, then $u\le h$.
\end{enumerate}
In particular $u$ is the supremum of $(\varphi_i)$ in the partially
ordered set $\PSH_X(\theta)$, and it does not depend on the choice of an
upper bound.
\end{proposition}

\begin{proof}
\cref{itm:ch45-lub:upper} is \Cref{thm:bj-final-envelope} together with the
fact that regularization increases a function.

\cref{itm:ch45-lub:usc}.  If $\varphi_i\le h$ for every $i$, then
$\sup_i\varphi_i\le h$ pointwise.  Since $h$ is upper semicontinuous and
$(\sup_i\varphi_i)^*$ is by definition the \emph{least} upper
semicontinuous majorant of $\sup_i\varphi_i$, we get $u\le h$.

\cref{itm:ch45-lub:least} is the special case $h=w$ of
\cref{itm:ch45-lub:usc}, a $\theta$-psh function being upper
semicontinuous by \Cref{def:bj-psh}\,(1).

The final sentence is the definition of a least upper bound; the
independence assertion holds because the construction of $u$ does not
mention the constant $C$ bounding the family, which enters only to
guarantee that $\sup_i\varphi_i\not\equiv+\infty$.
\end{proof}

\begin{remark}[Why the least upper bound is not the naive one]
\label{rem:ch45-lub-not-naive}
It is worth insisting that \cref{itm:ch45-lub:least} is not a formality.  In
a partially ordered set the least upper bound of a family need not exist,
and here the candidate --- the pointwise supremum --- is in general not a
member of the set at all.  The content of \Cref{prop:ch45-lub} is that
$\PSH_X(\theta)$ nevertheless has all the suprema one could want, and that
the supremum is computed by the concrete recipe ``take the pointwise
supremum, then regularize''.  Neither half of that recipe is available
before \Cref{thm:bj-final-envelope}: the first half because the result is
not psh, the second because regularization could a priori leave the class.
\end{remark}

\subsection{Compactness}
\label{sec:ch45-compactness}

\begin{corollary}
\label{cor:ch45-compactness}
Let $X$ be integral projective unibranch over an algebraically closed
trivially valued field, and let $\theta\in N^1(X)_{\mathbf R}$ be
arbitrary.  Then the normalized space
\[
 \PSH_X^{\sup}(\theta)
 =\{\phi\in\PSH_X(\theta)\ :\ \textstyle\sup_{X^{\An}}\phi=0\}
\]
of \eqref{eq:normalized-psh-space} is compact in the topology of pointwise
convergence on $X^{\divis}$.
\end{corollary}

\begin{proof}
By \Cref{thm:bj-final-envelope} the class $\theta$ has the envelope
property, and by \Cref{thm:bj-envelope-equivalences} the envelope property
is equivalent to the asserted compactness.  Note that
\eqref{eq:normalized-psh-space} makes sense for an arbitrary real class,
since $\sup_{X^{\An}}\phi$ is finite for every $\theta$-psh $\phi$;
and that the assertion is vacuously true when $\PSH_X(\theta)=\emptyset$,
which by \Cref{lem:ch34-psh-nonempty} happens in particular when $\theta$
is not pseudoeffective.  The converse is not needed here.
\end{proof}

\begin{remark}[Compactness is exactly as strong as the envelope property]
\label{rem:ch45-compactness-equivalent}
The implication in \Cref{cor:ch45-compactness} is reversible:
\Cref{thm:bj-envelope-equivalences} is an equivalence, so compactness of
$\PSH^{\sup}_X(\theta)$ implies $\operatorname{EP}_X(\theta)$.  Combined
with \Cref{thm:ep-implies-unibranch} this shows that on a variety that is
\emph{not} unibranch, $\PSH^{\sup}_X(\omega)$ fails to be compact for every
ample $\omega$.  The unibranch hypothesis in
\Cref{cor:ch45-compactness} is therefore not removable, and compactness of
the normalized psh space is not a soft consequence of the definitions: it is
the envelope theorem in another notation.  For nef classes the statement is
the conditional one recorded in \Cref{warn:ch34-compactness-conditional},
whose proof is the Tychonoff argument of \cite[Theorem~5.11]{BJ22} run on
the divisorial coordinates once the envelope property is available; what the
present corollary adds is that the hypothesis under which that argument
closes is now a theorem for unibranch $X$, and that the conclusion holds
for every real class, nef or not.  The unconditional form is false, as
\Cref{warn:ch34-compactness-conditional} explains.
\end{remark}

\subsection{Nets are not a formality}
\label{sec:ch45-nets-formality}

The formulations above are stated for nets because the directed set of
finite subsets of an arbitrary index set need not have a countable cofinal
subset; this is \Cref{prop:ch36-no-sequences}.  That obstruction concerns the
canonical finite-maximum reduction.  It does not exclude a separate
sequential approximation theorem for a particular family or function, and
none is assumed below.

\begin{warning}{No sequential reduction is used in this chapter}
\label{warn:ch45-nets}
The following four statements are formulated and proved for nets:
\begin{enumerate}
\item the approximation clause \Cref{def:bj-psh}\,(3) defining
      $\PSH_X(\theta)$;
\item the closure clauses \Cref{thm:bj-basic-closure}\,(4) and~(5);
\item the increasing-net criterion
      \Cref{thm:bj-final-increasing-nets};
\item the compactness statement \Cref{cor:ch45-compactness}, whose proof
      passes to a \emph{subnet} in the sense of the front matter, that is,
      along an order-preserving cofinal indexing map.
\end{enumerate}
Nothing in this chapter replaces those nets by sequences.  Wherever a
sequence appears below --- for instance in
\Cref{thm:bj-final-decreasing-obstacles} --- it is because the data are
indexed by $\mathbf N$ to begin with, never because a net has been
replaced by a sequence.
\end{warning}

The reader who wants a concrete instance of the failure of countable
cofinality may take $X$ a curve over an uncountable algebraically closed
field: already the family $(\varphi_a)_{a\in(0,1]}$ of
\Cref{ex:ch45-raw-sup-not-usc} is indexed by an uncountable totally ordered
set, and its net of finite maxima is indexed by the finite subsets of
$(0,1]$.  If a sequence of finite subsets were cofinal there, every
singleton would be contained in one of them; their countable union would
therefore equal $(0,1]$, which is impossible.  In this particular example a
cofinal
\emph{sequence} does exist in the underlying totally ordered set
$\bigl((0,1],\ge\bigr)$, namely $a=1/n$; the point is that no such
accident is available in general, and none is used.

\section{Perturbation of the class}
\label{sec:ch45-perturbation}

\Cref{thm:bj-final-envelope} is stated for an arbitrary real class, and the
proof in \Cref{sec:ch45-proof} treats an arbitrary real class directly.  It
is nevertheless useful to record the reduction that allows a proof to be run
in rational classes only, both because it is the reduction used inside
\Cref{chap:bj-smooth-characteristic} and because it is the form in which the
supplied manuscript quotes the perturbation lemma
\cite[Lemma~5.9]{BJ22}.

\subsection{The perturbation lemma}
\label{sec:ch45-perturbation-lemma}

The lemma has been proved once, in \Cref{chap:bj-smooth-characteristic}, and
we restate it here in the language of the present chapter because it will be
used again in \Cref{sec:ch45-obstacles} and in
\Crefrange{chap:bj-energy}{chap:bj-stability}.

Its content is the following asymmetry.  If $\theta'\ge\theta$ in the
sense that $\theta'-\theta$ is nef, then
\[
 \PSH_X(\theta)\subseteq\PSH_X(\theta'),
\]
because $0\in\PSH_X(\theta'-\theta)$ by \Cref{lem:ch34-psh-nonempty} and
sums of psh functions add classes, by \Cref{thm:bj-basic-closure}\,(3).  So
a family in the smaller class is a family in the larger one, with the
\emph{same} regularized supremum; if the larger class has the envelope
property, that supremum is $\theta'$-psh.  Letting $\theta'$ run through
a net converging to $\theta$ and using the closedness of
$\{\eta:\phi\in\PSH_X(\eta)\}$ --- which is
\Cref{thm:bj-basic-closure}\,(5) --- returns the supremum to
$\PSH_X(\theta)$.  This is \Cref{lem:ch40-ep-nef-limit}, and the
orientation of the inequality is essential: it is the inclusion
$\PSH_X(\theta)\subseteq\PSH_X(\theta_j)$, not its reverse, that makes the
argument run.

\begin{proposition}
\label{prop:ch45-reduction-big-rational}
Let $X$ be an integral projective $k$-variety and suppose that
\[
 \operatorname{EP}_X(c_1(L))\quad\text{holds for every big }
 \mathbf Q\text{-line bundle }L\text{ on }X .
\]
Then $\operatorname{EP}_X(\theta)$ holds for every
$\theta\in N^1(X)_{\mathbf R}$.
\end{proposition}

\begin{proof}
Let $\theta\in N^1(X)_{\mathbf R}$.  If $\theta$ is not
pseudoeffective, then $\PSH_X(\theta)=\emptyset$ by
\Cref{lem:ch34-psh-nonempty}, there is no nonempty family, and
$\operatorname{EP}_X(\theta)$ holds vacuously.

Assume therefore $\theta\in\operatorname{Psef}(X)$, and fix a norm on the
finite-dimensional real vector space $N^1(X)_{\mathbf R}$.  For each
$j\in\mathbf Z_{>0}$, \Cref{lem:ch40-rational-approximation} with
$\delta=1/j$ provides a $\mathbf Q$-line bundle $L_j$ with
\[
 \lVert c_1(L_j)-\theta\rVert<1/j,
 \qquad
 c_1(L_j)-\theta\ \text{ample},
 \qquad
 c_1(L_j)\ \text{big}.
\]
Set $\theta_j:=c_1(L_j)$.  By hypothesis $\operatorname{EP}_X(\theta_j)$
holds, since $L_j$ is a big $\mathbf Q$-line bundle.  Moreover
$\theta_j-\theta$ is ample, hence nef, and $\theta_j\to\theta$.
\Cref{lem:ch40-ep-nef-limit} now gives $\operatorname{EP}_X(\theta)$.
\end{proof}

\begin{corollary}
\label{cor:ch45-big-rational-suffices}
For an integral projective unibranch $X$, the assertion of
\Cref{thm:bj-final-envelope} for all $\theta\in N^1(X)_{\mathbf R}$ is
equivalent to its assertion for all classes of big $\mathbf Q$-line
bundles.
\end{corollary}

\begin{proof}
One direction is trivial and the other is
\Cref{prop:ch45-reduction-big-rational}.
\end{proof}

\begin{remark}[Why ample classes alone would not suffice]
\label{rem:ch45-ample-insufficient}
It is tempting to try to reduce further, to ample classes.  The attempt
fails at the first step, and the reason is worth recording, since the same
trap occurs in the energy chapters.  Given a pseudoeffective $\theta$ and
an ample $\alpha$, the perturbed classes $\theta+\varepsilon\alpha$ do
converge to $\theta$ with nef --- indeed ample --- difference, which is what
\Cref{lem:ch40-ep-nef-limit} requires; but $\theta+\varepsilon\alpha$ is
ample only when $\theta$ is nef, and a pseudoeffective class need not be
nef.  What \Cref{lem:ch40-rational-approximation} does provide is bigness of
$\theta+\varepsilon\alpha$, by an asymptotic Riemann--Roch argument, and
bigness is exactly the positivity that the obstacle machinery of
\Cref{chap:bj-smooth-characteristic} needs: a big class has enough sections
in large multiples to produce Fubini--Study competitors, while an ample class
would be a strictly stronger and unavailable hypothesis.  So the correct
reduction is to \emph{big} rational classes, and
\Cref{prop:ch45-reduction-big-rational} is sharp in that respect.
\end{remark}

\begin{remark}[The two orientations of a class perturbation]
\label{rem:ch45-two-orientations}
\Cref{prop:ch38-drift-ledger}\,\cref{itm:ch38-drift:orient} and
\Cref{warn:ch38-drift-not-absorbed} distinguish two uses of a small class
perturbation, and both occur in the envelope proof.  In the reduction just
performed, the perturbed classes \emph{dominate} the target and the
functions are lifted from $\PSH_X(\theta)$ into $\PSH_X(\theta_j)$; the
relevant closure clause is \Cref{thm:bj-basic-closure}\,(3).  In the
approximation arguments of
\Cref{chap:bj-smooth-characteristic,chap:bj-subintegral-normalization} the
perturbed classes also dominate the target, but the functions are pushed
\emph{down} to it along a decreasing net; the relevant clauses are
\Cref{thm:bj-basic-closure}\,(4) and~(5).  The two are not interchangeable,
and \Cref{rem:ch40-reduction-orientation} keeps them apart at the point where
both occur within a page of each other.
\end{remark}

\begin{remark}[Real classes again]
\label{rem:ch45-real-again}
\Cref{rem:ch37-real-classes} explained why a proof valid only for rational
classes would not suffice: the passage to the limit would require knowing in
advance that the limit of the envelopes is the envelope of the limit.
\Cref{prop:ch45-reduction-big-rational} is not a counterexample to that
remark, and it is worth seeing why.  It does not compare envelopes in
different classes at all; it observes that one and the same function ---
the regularized supremum of one and the same family --- lies in
$\PSH_X(\theta_j)$ for every $j$, and then applies the closedness of the
set of classes for which a \emph{fixed} function is psh.  That closedness is
\Cref{thm:bj-basic-closure}\,(5), which is a consequence of the
varying-class convention in \Cref{def:bj-psh}\,(3), and it is available
precisely because the definition of a psh function was set up with classes
allowed to vary.
\end{remark}

\section{Obstacle envelopes in arbitrary real classes}
\label{sec:ch45-obstacles}

\Cref{chap:bj-envelope-history} introduced the obstacle envelope
$P_\theta$ of \Cref{def:bj-obstacle-envelope} and immediately attached a
warning to it: before the envelope theorem, $P_\theta$ is a map from
obstacles to \emph{functions}, not to $\PSH_X(\theta)$, and writing
$P_\theta:\PSH_X(\theta)\to\PSH_X(\theta)$ or composing $P_\theta$ with
itself presupposes exactly what is at issue.  This section lifts that
warning.

\begin{important}{The obstacle envelope is now an operator}
\label{imp:ch45-obstacle-operator}
Let $X$ be integral projective unibranch and
$\theta\in N^1(X)_{\mathbf R}$.  For every upper semicontinuous
$f:X^{\An}\to\mathbf R\cup\{-\infty\}$ that is bounded above and admits at
least one $\theta$-psh competitor, the function $P_\theta(f)$ of
\Cref{def:bj-obstacle-envelope} is $\theta$-psh, is already upper
semicontinuous without regularization, and is the greatest $\theta$-psh
function dominated by $f$.  From this point on, $P_\theta$ may be used as
an operator
\[
 P_\theta:\ \{\text{admissible obstacles}\}\longrightarrow\PSH_X(\theta),
\]
and in particular as a projection $P_\theta:\PSH_X(\theta)\to\PSH_X(\theta)$,
which is the identity.  The warning of \Cref{chap:bj-envelope-history} is
discharged; the hypotheses ``upper semicontinuous, bounded above, at least
one competitor'' are not, and are carried in every statement below.
\end{important}

\subsection{The two definitions agree}
\label{sec:ch45-two-definitions}

\Cref{chap:bj-envelope-history} defined $P_\theta(h)$ as a raw supremum,
without regularization, precisely because at that stage the regularization
could not be shown to remain psh.  The definition used in the present
chapter carries a star.  The following lemma shows that for admissible
obstacles the two agree, so that the notation $P_\theta$ is unambiguous.

\begin{definition}[Admissible obstacle]
\label{def:ch45-admissible-obstacle}
Let $\theta\in N^1(X)_{\mathbf R}$.  A function
$f:X^{\An}\to\mathbf R\cup\{-\infty\}$ is an \emph{admissible obstacle for
$\theta$} if
\begin{enumerate}
\item $f$ is upper semicontinuous;
\item $f$ is bounded above by a real constant on $X^{\An}$;
\item the competitor family
      $\mathcal F_\theta(f):=\{\psi\in\PSH_X(\theta)\mid\psi\le f\}$
      is nonempty.
\end{enumerate}
\end{definition}

\begin{lemma}
\label{lem:ch45-raw-equals-regularized}
Let $X$ be integral projective unibranch, $\theta\in N^1(X)_{\mathbf R}$,
and let $f$ be an admissible obstacle for $\theta$.  Then the raw
supremum
\[
 S:=\sup\mathcal F_\theta(f)
\]
is upper semicontinuous, belongs to $\PSH_X(\theta)$, and satisfies
$S=S^*$.  Consequently
\Cref{def:bj-obstacle-envelope} and
\Cref{def:bj-final-obstacle-envelope} define the same function, and the
symbol $P_\theta(f)$ is unambiguous.
\end{lemma}

\begin{proof}
The family $\mathcal F_\theta(f)$ is nonempty by admissibility and bounded above
by the constant bounding $f$.  By \Cref{thm:bj-final-envelope},
$S^*\in\PSH_X(\theta)$.

We claim $S^*\le f$.  Indeed $S\le f$ pointwise, since every member of
$\mathcal F_\theta(f)$ is $\le f$; and $f$ is upper semicontinuous, hence an
upper semicontinuous majorant of $S$; since $S^*$ is the \emph{least}
such majorant, $S^*\le f$.

Therefore $S^*\in\mathcal F_\theta(f)$, so $S^*\le\sup\mathcal F_\theta(f)=S$.  The
reverse inequality $S\le S^*$ holds by definition of the regularization.
Hence $S=S^*$, which is upper semicontinuous and $\theta$-psh.

The last assertion is now immediate: \Cref{def:bj-obstacle-envelope} gives
$S$ and \Cref{def:bj-final-obstacle-envelope} gives $S^*$, and these are
equal.
\end{proof}

\begin{remark}[Where admissibility is used, and what fails without it]
\label{rem:ch45-admissibility}
Each of the three clauses of \Cref{def:ch45-admissible-obstacle} is used
exactly once in the proof of \Cref{lem:ch45-raw-equals-regularized}, and
each fails in a different way if dropped.

Without nonemptiness of $\mathcal F_\theta(f)$ the supremum of the empty family
is the constant $-\infty$, which is not a psh function by
\Cref{def:bj-psh}\,(2); this is the first alternative of the trichotomy
\eqref{eq:envelope-trichotomy}.

Without an upper bound the competitors may have suprema tending to
$+\infty$, and the regularized supremum is then identically $+\infty$ by
the Alexander--Taylor argument recorded in the proof of
\Cref{thm:general-obstacle-criterion}; this is the second alternative.

Without upper semicontinuity of $f$ the inequality $S^*\le f$ may fail,
and then $S^*$ is psh but is not a competitor, so the raw and regularized
envelopes genuinely differ.  Here is the smallest instance.  Take
$X=\mathbf P^1$, $\theta=c_1(\cO(1))$, a closed point $p$, and let
\[
 f(v):=\begin{cases}-1,& v=v_{p,\triv},\\ 0,&\text{otherwise.}\end{cases}
\]
Then $f$ is bounded above and fails upper semicontinuity at
$v_{p,\triv}$, since $f(t\ord_p)=0$ for every finite $t$ while
$t\ord_p\to v_{p,\triv}$.  Its competitor family contains the functions
$\max\{-at,-1\}$ of \Cref{ex:ch45-raw-sup-not-usc}, truncated at the level
$-1$: explicitly, for $a\in(0,1]$ let $\psi_a$ be $0$ off the ray
through $\ord_p$ and $\psi_a(t\ord_p)=\max\{-at,-1\}$, which is
$\theta$-psh by the curve criterion and satisfies
$\psi_a(v_{p,\triv})=-1\le f(v_{p,\triv})$ and $\psi_a\le0=f$ elsewhere.
Hence
\[
 S=\sup\mathcal F_\theta(f)=\begin{cases}-1,& v=v_{p,\triv},\\ 0,&\text{otherwise,}\end{cases}
\]
the value $0$ off $v_{p,\triv}$ arising as $a\downarrow0$; so
$S=f$, $S$ is not upper semicontinuous, and $S^*\equiv0$.  In
particular $S^*\not\le f$, the raw and starred definitions disagree, and
the starred one produces a function that is not a competitor at all.  Upper
semicontinuity of the obstacle is thus exactly the hypothesis that makes
\Cref{lem:ch45-raw-equals-regularized} true, and it is a hypothesis on the
obstacle, not on the class or on the variety.
\end{remark}

\subsection{The envelope is psh}
\label{sec:ch45-obstacle-psh}

\begin{definition}[Obstacle envelope]
\label{def:bj-final-obstacle-envelope}
For an upper semicontinuous function
$f:X^{\An}\to\mathbf R\cup\{-\infty\}$, bounded above, set
\[
 P_\theta(f):=
 \left(
   \sup\{\varphi\in\PSH_X(\theta)\mid \varphi\le f\}
 \right)^*
\tag{45.8}\label{eq:bj-final-obstacle-envelope}
\]
whenever the competitor family is nonempty.
\end{definition}

\begin{theorem}\label{thm:bj-final-obstacle-psh}
If $X$ is unibranch, the function $P_\theta(f)$ in
\eqref{eq:bj-final-obstacle-envelope} belongs to $\PSH_X(\theta)$.
It satisfies $P_\theta(f)\le f$ and is the largest
$\theta$-psh function dominated by $f$.
\end{theorem}

\begin{proof}
The competitor family is nonempty by hypothesis and is bounded above by the
finite global upper bound for $f$.
\Cref{thm:bj-final-envelope} makes its regularized supremum psh.
The raw supremum is at most $f$, and $f$ is upper semicontinuous, so its
upper semicontinuous regularization is still at most $f$.
Every competitor is at most the regularized supremum.
This proves maximality.

In the language of \Cref{def:ch45-admissible-obstacle}, the hypotheses of
the theorem say exactly that $f$ is an admissible obstacle for $\theta$,
and the conclusion is that $\mathcal F_\theta(f)$ has a greatest element, namely
$P_\theta(f)$.  Existence of a greatest element is a strictly stronger
assertion than existence of a least upper bound, and it is worth separating
the two: \Cref{prop:ch45-lub} produces the least upper bound of a family
inside $\PSH_X(\theta)$, while the present theorem produces an element of
the family $\mathcal F_\theta(f)$ dominating all the others.  The passage from
one to the other is \Cref{lem:ch45-raw-equals-regularized}: the least upper
bound of $\mathcal F_\theta(f)$ happens to be a member of $\mathcal F_\theta(f)$
because $f$ is upper semicontinuous.
\end{proof}

\subsection{The calculus of \texorpdfstring{$P_\theta$}{P theta}}
\label{sec:ch45-obstacle-calculus}

\begin{proposition}\label{prop:bj-final-obstacle-calculus}
Let $X$ be integral projective unibranch and
$\theta\in N^1(X)_{\mathbf R}$, and let $f,g$ be admissible obstacles
for $\theta$.  Then:
\begin{enumerate}
\item if $f\le g$, then $P_\theta(f)\le P_\theta(g)$;
\item for $c\in\mathbf R$,
      $P_\theta(f+c)=P_\theta(f)+c$;
\item $P_\theta(P_\theta(f))=P_\theta(f)$;
\item the construction is valid for every real
      $\theta\in N^1(X)_{\mathbf R}$, not only rational line-bundle classes.
\end{enumerate}
\par
\end{proposition}

\begin{proof}
The first assertion follows by inclusion of competitor families.
Translation by $c$ gives a bijection between the competitors for $f$ and
those for $f+c$.
The theorem says that $P_\theta(f)$ is itself a competitor for its own
obstacle, which proves idempotence.
The last assertion is the class range of
\Cref{thm:bj-final-envelope}.

We write the four arguments out, since \Cref{lem:obstacle-envelope-formal}
proved the same formal identities for the raw supremum and it is the
upgrade, not the identities, that is at stake.

\emph{(1) Monotonicity.}  If $f\le g$ then every competitor for $f$ is a
competitor for $g$, so $\mathcal F_\theta(f)\subseteq\mathcal F_\theta(g)$ and hence
$\sup\mathcal F_\theta(f)\le\sup\mathcal F_\theta(g)$; regularization is order
preserving, so $P_\theta(f)\le P_\theta(g)$.  Note that the hypothesis
that \emph{both} families be nonempty is not needed here: nonemptiness for
$f$ implies it for $g$.

\emph{(2) Constants.}  The map $\psi\mapsto\psi+c$ is a bijection of
$\PSH_X(\theta)$ onto itself, by \Cref{thm:bj-basic-closure}\,(1), and it
carries $\mathcal F_\theta(f)$ bijectively onto $\mathcal F_\theta(f+c)$.  Hence
$\sup\mathcal F_\theta(f+c)=\sup\mathcal F_\theta(f)+c$; and
$(\cdot+c)$ commutes with upper semicontinuous regularization, because
translation by a real constant is a homeomorphism of $\Rbar$ fixing
$\pm\infty$ and preserving the order.  Thus
$P_\theta(f+c)=P_\theta(f)+c$.

\emph{(3) Idempotence.}  Write $u:=P_\theta(f)$.  By
\Cref{thm:bj-final-obstacle-psh}, $u\in\PSH_X(\theta)$ and $u\le f$, so
$u\in\mathcal F_\theta(u)$; on the other hand every $\psi\in\mathcal F_\theta(u)$
satisfies $\psi\le u$.  Hence $u$ is the greatest element of
$\mathcal F_\theta(u)$, that is, $\sup\mathcal F_\theta(u)=u$, and $u$ is upper
semicontinuous, so $P_\theta(u)=u$.  This is the assertion, and it is
exactly the place where \Cref{lem:obstacle-envelope-formal}\,(3) had to
hedge: composing $P_\theta$ with itself requires knowing that the inner
value is again an obstacle of the admissible kind, which is
\Cref{thm:bj-final-obstacle-psh}.

\emph{(4) Class range.}  \Cref{thm:bj-final-envelope} is proved for an
arbitrary $\theta\in N^1(X)_{\mathbf R}$, and no step of (1)--(3) used
rationality, integrality or positivity of the class.  The only positivity
that enters is hidden in the nonemptiness hypothesis: by
\Cref{lem:ch34-psh-nonempty}, $\mathcal F_\theta(f)\ne\emptyset$ forces
$\theta$ to be pseudoeffective, and forces $\theta$ to be nef as soon as
some competitor is bounded.
\end{proof}

The calculus continues.  The next proposition collects the identities that
are used in \Crefrange{chap:bj-energy}{chap:bj-stability}; all of them are
false, or meaningless, before \Cref{thm:bj-final-obstacle-psh}.

\begin{proposition}
\label{prop:ch45-obstacle-calculus-2}
Let $X$ be integral projective unibranch,
$\theta\in N^1(X)_{\mathbf R}$, and let $f,g$ be admissible obstacles
for $\theta$.  Then:
\begin{enumerate}
\item\label{itm:ch45-P:projection}
      $P_\theta(u)=u$ for every $u\in\PSH_X(\theta)$ that is bounded
      above --- which is every $u\in\PSH_X(\theta)$; so $P_\theta$
      restricts to the identity on $\PSH_X(\theta)$;
\item\label{itm:ch45-P:lipschitz}
      if $c\ge0$ and $f\le g+c$, then
      $P_\theta(f)\le P_\theta(g)+c$; consequently, if $\theta$ is nef
      and $f,g$ are real valued and bounded, then $P_\theta(f)$ and
      $P_\theta(g)$ are bounded and
      $\lVert P_\theta(f)-P_\theta(g)\rVert_\infty
        \le\lVert f-g\rVert_\infty$;
\item\label{itm:ch45-P:divisorial}
      if $f$ and $g$ agree on $X^{\divis}$, then
      $P_\theta(f)=P_\theta(g)$ on all of $X^{\An}$;
\item\label{itm:ch45-P:max}
      $\max\{f,g\}$ is an admissible obstacle, and
      \[
       \max\{P_\theta(f),P_\theta(g)\}
       \le P_\theta(\max\{f,g\}).
      \]
      If $\min\{f,g\}$ is an admissible obstacle, then
      \[
       P_\theta(\min\{f,g\})
       \le\min\{P_\theta(f),P_\theta(g)\}.
      \]
\item\label{itm:ch45-P:normalization}
      $\nu^*P_\theta(f)=P_{\nu^*\theta}(\nu^*f)$, where
      $\nu:X^\nu\to X$ is the normalization;
\item\label{itm:ch45-P:envelope-of-family}
      if $(\varphi_i)_{i\in I}\subset\PSH_X(\theta)$ is nonempty and
      bounded above, then
      $P_\theta\bigl((\sup_i\varphi_i)^*\bigr)=(\sup_i\varphi_i)^*$.
\end{enumerate}
\end{proposition}

\begin{proof}
\cref{itm:ch45-P:projection}.  Every $u\in\PSH_X(\theta)$ is bounded above,
by the finiteness of $\sup_{X^{\An}}u$ recorded after
\eqref{eq:normalized-psh-space}, and is upper semicontinuous by
\Cref{def:bj-psh}\,(1); and $u\in\mathcal F_\theta(u)$, so $u$ is an admissible
obstacle.  Since $u$ is the greatest element of $\mathcal F_\theta(u)$, we get
$P_\theta(u)=u$.  This is \Cref{prop:bj-final-obstacle-calculus}\,(3)
applied to $f=u$, and it identifies $P_\theta$ as a projection.

\cref{itm:ch45-P:lipschitz}.  From $f\le g+c$ and
\Cref{prop:bj-final-obstacle-calculus}\,(1),(2) we get
$P_\theta(f)\le P_\theta(g+c)=P_\theta(g)+c$.  Suppose now that
$\theta$ is nef and that $f,g$ are real valued and bounded.  By
\Cref{lem:ch34-psh-nonempty} the constants are $\theta$-psh, so
$\inf f$ is a competitor for $f$ and $\inf f\le P_\theta(f)\le\sup f$;
in particular $P_\theta(f)$ and $P_\theta(g)$ are real valued and
bounded, and their difference is defined everywhere.  Apply the first
assertion with $c=\lVert f-g\rVert_\infty$, and then again with $f$ and
$g$ interchanged.

\cref{itm:ch45-P:divisorial}.  Let $\psi\in\PSH_X(\theta)$.  By the
comparison interface \cite[Theorem~4.22]{BJ22}, recorded in
\Cref{chap:bj-envelope-statement}, and since $f$ is upper semicontinuous,
$\psi\le f$ on $X^{\An}$ if and only if $\psi\le f$ on
$X^{\divis}$.  The same holds for $g$.  As $f=g$ on $X^{\divis}$,
the two conditions coincide, so $\mathcal F_\theta(f)=\mathcal F_\theta(g)$ and the two
envelopes are equal.  This is a genuinely useful statement: it says that an
obstacle is only ever felt on $X^{\divis}$, which is where obstacles are
computed in practice.

\cref{itm:ch45-P:max}.  The function $\max\{f,g\}$ is upper
semicontinuous, is bounded above by the maximum of the two bounds, and
admits a competitor, namely any competitor for $f$; so it is admissible.
Both inequalities are then instances of
\Cref{prop:bj-final-obstacle-calculus}\,(1), since $f\le\max\{f,g\}\ge g$
and $\min\{f,g\}\le f,g$.  For the first inequality one may argue
directly: $\max\{P_\theta(f),P_\theta(g)\}$ is $\theta$-psh by
\Cref{thm:bj-basic-closure}\,(2) and is $\le\max\{f,g\}$, hence is a
competitor.  Admissibility of $\min\{f,g\}$ is not automatic --- its
competitor family may be empty --- and is therefore hypothesized.

\cref{itm:ch45-P:normalization}.  By
\Cref{lem:ch45-normalization-data}\,\cref{itm:ch45-nd:uh} the map
$\nu^{\An}$ is a homeomorphism, so $\nu^*f$ is upper semicontinuous and
bounded above by the same constant.  By
\Cref{cor:ch44-unibranch-normalization-reflection}\,\cref{itm:ch44-unib:refl}
the map $\nu^*$ is an order-preserving and order-reflecting bijection
$\PSH_X(\theta)\to\PSH_{X^\nu}(\nu^*\theta)$.  Hence
$\nu^*$ carries $\mathcal F_\theta(f)$ bijectively onto
$\mathcal F_{\nu^*\theta}(\nu^*f)$, and greatest elements correspond.  Since
$P_\theta(f)$ is the greatest element of $\mathcal F_\theta(f)$ by
\Cref{thm:bj-final-obstacle-psh}, its image $\nu^*P_\theta(f)$ is the
greatest element of $\mathcal F_{\nu^*\theta}(\nu^*f)$, which is
$P_{\nu^*\theta}(\nu^*f)$.

\cref{itm:ch45-P:envelope-of-family}.  By \Cref{thm:bj-final-envelope} the
function $u:=(\sup_i\varphi_i)^*$ lies in $\PSH_X(\theta)$; apply
\cref{itm:ch45-P:projection}.
\end{proof}

\begin{corollary}
\label{cor:ch45-extremal}
Let $X$ be integral projective unibranch and let
$\theta\in N^1(X)_{\mathbf R}$ with $\PSH_X(\theta)\ne\emptyset$.  Then
the constant $0$ is an admissible obstacle for $\theta$, the function
\[
 V_\theta:=P_\theta(0)\in\PSH_X(\theta)
\]
is the greatest $\theta$-psh function that is $\le0$, and
\[
 V_\theta\equiv0
 \quad\Longleftrightarrow\quad
 \theta\in\Nef(X).
\]
Moreover $P_\theta(c)=V_\theta+c$ for every $c\in\mathbf R$.
\end{corollary}

\begin{proof}
Admissibility: the constant $0$ is continuous, hence upper semicontinuous,
and bounded; and if $\psi\in\PSH_X(\theta)$ is arbitrary, then
$\psi-\sup_{X^{\An}}\psi\in\PSH_X(\theta)$ is $\le0$ by
\eqref{eq:normalized-psh-space} and
\Cref{thm:bj-basic-closure}\,(1), so $\mathcal F_\theta(0)\ne\emptyset$.
\Cref{thm:bj-final-obstacle-psh} then gives the first two assertions.

If $\theta$ is nef, the constant $0$ is $\theta$-psh by
\Cref{lem:ch34-psh-nonempty}, hence is itself the greatest element of
$\mathcal F_\theta(0)$, so $V_\theta=0$.  Conversely, if $V_\theta=0$ then
$0\in\PSH_X(\theta)$ is a bounded member, and
\eqref{eq:psh-existence-cones} gives $\theta\in\Nef(X)$.

The last identity is \Cref{prop:bj-final-obstacle-calculus}\,(2).
\end{proof}

\begin{remark}[\texorpdfstring{$V_\theta$}{V theta} is where the two positivity thresholds separate]
\label{rem:ch45-two-thresholds}
\Cref{lem:ch34-psh-nonempty} recorded two different thresholds:
$\PSH_X(\theta)\ne\emptyset$ needs $\theta$ pseudoeffective, while a
\emph{bounded} member needs $\theta$ nef.  \Cref{cor:ch45-extremal} makes
the gap between them into a single function whenever
$\PSH_X(\theta)\ne\emptyset$.  If in addition $\theta$ is not nef, then
$V_\theta$ is a well-defined $\theta$-psh function that is $\le0$, is
not identically $0$, and is unbounded below;
it is the non-Archimedean counterpart of the equilibrium potential of a
pseudoeffective class.  Its existence is a consequence of the envelope
theorem and is not available before this chapter.
\end{remark}

\subsection{The trichotomy, resolved}
\label{sec:ch45-trichotomy}

\begin{proposition}
\label{prop:ch45-trichotomy}
Let $X$ be integral projective unibranch, $\theta\in N^1(X)_{\mathbf R}$
with $\PSH_X(\theta)\ne\emptyset$, and let
$h:X^{\An}\to\Rbar$ be arbitrary.  Then exactly one of the following holds:
\begin{enumerate}
\item $\mathcal F_\theta(h)=\emptyset$, and $P_\theta(h)\equiv-\infty$;
\item $\mathcal F_\theta(h)\ne\emptyset$ and is not bounded above by any real
      constant, and $P_\theta(h)^*\equiv+\infty$;
\item $\mathcal F_\theta(h)\ne\emptyset$ and is bounded above, and
      $P_\theta(h)^*\in\PSH_X(\theta)$.
\end{enumerate}
If moreover $h$ is upper semicontinuous and bounded above and
$\mathcal F_\theta(h)\ne\emptyset$, then the third alternative holds and
$P_\theta(h)=P_\theta(h)^*$.
\end{proposition}

\begin{proof}
The trichotomy is \Cref{thm:general-obstacle-criterion}\,(1), which is an
equivalence with $\operatorname{EP}_X(\theta)$; the latter holds by
\Cref{thm:bj-final-envelope}.  The three alternatives as stated are mutually
exclusive and exhaustive by inspection of the two conditions
``$\mathcal F_\theta(h)=\emptyset$'' and ``$\mathcal F_\theta(h)$ bounded above'', and
the third case is \Cref{thm:bj-final-envelope} applied to the family
$\mathcal F_\theta(h)$.  In the second case the members of $\mathcal F_\theta(h)$ have
suprema tending to $+\infty$, and the Alexander--Taylor bound quoted in
the proof of \Cref{thm:general-obstacle-criterion} forces the regularization
to be $+\infty$ at every divisorial point, hence identically $+\infty$
by density.

For the last sentence: if $h$ is bounded above by $C$, then so is
$\mathcal F_\theta(h)$, which excludes the second alternative; and
$\mathcal F_\theta(h)\ne\emptyset$ is the remaining hypothesis, so the third
alternative holds.  \Cref{lem:ch45-raw-equals-regularized} then gives
$P_\theta(h)=P_\theta(h)^*$.
\end{proof}

\begin{corollary}
\label{cor:ch45-continuity-envelopes}
Let $X$ be integral projective unibranch and $\omega\in\Amp(X)$.  Then
for every $h\in C^0(X^{\An})$,
\[
 P_\omega(h)\in C^0(X^{\An})\cap\PSH_X(\omega),
\]
and $h\mapsto P_\omega(h)$ is a $1$-Lipschitz retraction of
$C^0(X^{\An})$ onto $C^0(X^{\An})\cap\PSH_X(\omega)$.
\end{corollary}

\begin{proof}
An ample class is nef, so $\PSH_X(\omega)$ contains the constants by
\Cref{lem:ch34-psh-nonempty}; a continuous $h$ on the compact space
$X^{\An}$ is bounded, so the constant $\inf h$ is a competitor and $h$
is an admissible obstacle.  \Cref{thm:bj-final-obstacle-psh} gives
$P_\omega(h)\in\PSH_X(\omega)$.  Continuity is the criterion
\eqref{eq:continuity-of-envelopes} recorded in
\Cref{chap:bj-envelope-history}, which for an ample class is equivalent to
$\operatorname{EP}_X(\omega)$; the latter holds by
\Cref{thm:bj-final-envelope}.  The Lipschitz property is
\Cref{prop:ch45-obstacle-calculus-2}\,\cref{itm:ch45-P:lipschitz}, and the
retraction property is
\Cref{prop:ch45-obstacle-calculus-2}\,\cref{itm:ch45-P:projection}
together with the fact that a continuous $\omega$-psh function is its own
envelope.
\end{proof}

\begin{warning}{Continuity is an ample-class phenomenon}
\label{warn:ch45-continuity-ample}
\Cref{cor:ch45-continuity-envelopes} is stated for an ample class and the
hypothesis is used, not decorative.  The equivalence
\eqref{eq:continuity-of-envelopes} is proved in
\Cref{chap:bj-envelope-history} for $\omega$ ample, because ampleness is
what supplies enough finite-valued Fubini--Study approximants.  If
$\PSH_X(\theta)\ne\emptyset$ and $\theta$ is not nef, then the envelope
of the continuous obstacle $0$ is $V_\theta$, which is unbounded below:
a bounded $\theta$-psh function would force $\theta\in\Nef(X)$ by
\Cref{lem:ch34-psh-nonempty}.  In particular $V_\theta\notin
C^0(X^{\An})$, a continuous function on a compact space being bounded.  It
does \emph{not} follow that $V_\theta$ fails to be finite valued: an upper
semicontinuous function on a compact space can be real valued and unbounded
below, as $x\mapsto-1/x$ on $(0,1]$, extended by $0$ at $0$,
already shows.  The envelope theorem is a
closure statement and supplies no regularity;
\Cref{rem:ch37-not-said} said so in advance, and
\Cref{ex:ch45-no-regularity} says it with a computation.
\end{warning}

\begin{remark}[Real class drift in the proof]
\label{rem:ch45-real-class-drift}
The final theorem is stated in one fixed real class, but its proof used GF
stages in rational classes converging to that class.
Class promotion in \Cref{chap:bj-universal-homeomorphisms},
determinant-norm normalization in \Cref{chap:bj-finite-descent},
and the error term in \eqref{eq:bj-char-zero-gf-approximants} are what make
that passage legitimate.
There is no hidden assumption that $\theta=c_1(L)$ for a rational line
bundle.
\end{remark}

\section{Monotone approximation}
\label{sec:ch45-monotone}

The envelope theorem is a closure statement.  This section converts it into
a calculus of monotone limits, which is the form in which
\Crefrange{chap:bj-energy}{chap:bj-stability} consume it.

\subsection{Decreasing obstacles}
\label{sec:ch45-decreasing-obstacles}

\begin{theorem}\label{thm:bj-final-decreasing-obstacles}
Let $X$ be integral projective unibranch and
$\theta\in N^1(X)_{\mathbf R}$.  Let $f_j\downarrow f$ be a decreasing
sequence of bounded-above upper semicontinuous obstacles.
Assume that the competitor family for $f$ is nonempty.
Then
\[
                         P_\theta(f_j)\downarrow P_\theta(f)
\tag{45.9}\label{eq:bj-final-decreasing-obstacles}
\]
pointwise on $X^{\An}$.
\end{theorem}

\begin{proof}
Monotonicity gives a decreasing sequence $P_\theta(f_j)$.
Its classes are all $\theta$.
A fixed competitor for $f$ lies below every term, so the limit is not
identically $-\infty$; PSH2 makes the limit $v$ a
$\theta$-psh function.
Since $P_\theta(f_j)\le f_j$, passage to the pointwise limit gives
$v\le f$.
Thus $v$ is a competitor for $f$, so $v\le P_\theta(f)$.
Conversely, monotonicity gives
$P_\theta(f)\le P_\theta(f_j)$ for all $j$.
Taking the limit yields the reverse inequality.

For later use we record the four checks that the argument performs, since
each of them is a hypothesis somewhere.  First, each $f_j$ is an
admissible obstacle in the sense of
\Cref{def:ch45-admissible-obstacle}: it is upper semicontinuous and bounded
above by hypothesis, and its competitor family contains
$\mathcal F_\theta(f)\ne\emptyset$ because $f\le f_j$.  Second, the sequence
$\bigl(P_\theta(f_j)\bigr)_j$ is decreasing by
\Cref{prop:bj-final-obstacle-calculus}\,(1) and lies in the fixed class
$\theta$, so the applicable closure clause is
\Cref{thm:bj-basic-closure}\,(4) rather than the varying-class
clause~(5).  Third, the exclusion of the identically $-\infty$
alternative in that clause is supplied by any fixed
$\psi\in\mathcal F_\theta(f)$, which satisfies $\psi\le P_\theta(f)\le
P_\theta(f_j)$ for every $j$ and is generically finite.  Fourth, the
passage to the limit in $P_\theta(f_j)\le f_j$ is a passage to the limit
in an inequality between decreasing families, which is legitimate in
$\Rbar$ with no subtraction.
\end{proof}

\begin{theorem}
\label{thm:ch45-decreasing-obstacle-nets}
Let $X$ be integral projective unibranch and
$\theta\in N^1(X)_{\mathbf R}$.  Let $(f_j)_{j\in J}$ be a downward
directed net of upper semicontinuous functions
$X^{\An}\to\mathbf R\cup\{-\infty\}$, bounded above by one real constant,
and put $f:=\inf_jf_j$.  If $\mathcal F_\theta(f)\ne\emptyset$, then $f$ is an
admissible obstacle and
\[
 \inf_{j\in J}P_\theta(f_j)=P_\theta(f)
\]
pointwise on $X^{\An}$.
\end{theorem}

\begin{proof}
The function $f$ is an infimum of upper semicontinuous functions, hence
upper semicontinuous, and is bounded above by the same constant; with the
hypothesis $\mathcal F_\theta(f)\ne\emptyset$ it is admissible.  Each $f_j$ is
admissible, since $f\le f_j$ gives
$\mathcal F_\theta(f)\subseteq\mathcal F_\theta(f_j)$.

Write $v:=\inf_jP_\theta(f_j)$.  The net $(P_\theta(f_j))_j$ is
downward directed in $\PSH_X(\theta)$ by
\Cref{prop:bj-final-obstacle-calculus}\,(1), so it is a decreasing net in
the sense of \Cref{thm:bj-basic-closure}\,(4).  Any fixed
$\psi\in\mathcal F_\theta(f)$ satisfies $\psi\le P_\theta(f_j)$ for all $j$,
so $v\ge\psi$ is generically finite, and
\Cref{thm:bj-basic-closure}\,(4) gives $v\in\PSH_X(\theta)$.

Since $P_\theta(f_j)\le f_j$ for each $j$, taking infima gives
$v\le f$; hence $v\in\mathcal F_\theta(f)$ and $v\le P_\theta(f)$.
Conversely $f\le f_j$ gives $P_\theta(f)\le P_\theta(f_j)$ for every
$j$, so $P_\theta(f)\le v$.
\end{proof}

\begin{remark}[Nets, and why \Cref{thm:bj-final-decreasing-obstacles} is
stated for sequences]
\label{rem:ch45-nets-vs-sequences-obstacles}
\Cref{thm:ch45-decreasing-obstacle-nets} contains
\Cref{thm:bj-final-decreasing-obstacles} as the special case
$J=\mathbf N$ with the reverse order.  The sequential form is retained
because the obstacles that arise in practice --- base-ideal stages
$\phi_D+m^{-1}\phi_{\mathfrak a_m}$, Frobenius stages, additive-power
blocks --- are indexed by $\mathbf N$ or by $\mathbf N$ under
divisibility, and stating those applications for nets would suggest that a
net has been chosen where none was needed.  This is the discipline of
\Cref{warn:ch45-nets}: nets where nets are forced, sequences where the data
are sequences.
\end{remark}

\subsection{Increasing obstacles}
\label{sec:ch45-increasing-obstacles}

\begin{theorem}
\label{thm:ch45-increasing-obstacles}
Let $X$ be integral projective unibranch and
$\theta\in N^1(X)_{\mathbf R}$.  Let $(h_j)_{j\in J}$ be an upward
directed net in $C^0(X^{\An})$ with $h:=\sup_jh_j$, and assume that
$\mathcal F_\theta(h_{j_0})\ne\emptyset$ for some $j_0$ and that
$\mathcal F_\theta(h)$ is bounded above by a real constant.  Then
\begin{enumerate}
\item\label{itm:ch45-inc:psh}
      $P_\theta(h_j)\in\PSH_X(\theta)$ for every $j\ge j_0$, and the net
      $(P_\theta(h_j))_{j\ge j_0}$ is increasing;
\item\label{itm:ch45-inc:sup}
      $\sup_{j}P_\theta(h_j)=\sup\mathcal F_\theta(h)$ pointwise on
      $X^{\An}$;
\item\label{itm:ch45-inc:star}
      $\bigl(\sup_jP_\theta(h_j)\bigr)^*
       =\bigl(\sup\mathcal F_\theta(h)\bigr)^*\in\PSH_X(\theta)$, and this
      function is the least upper bound of $\{P_\theta(h_j)\}_j$ in
      $\PSH_X(\theta)$.
\end{enumerate}
\end{theorem}

\begin{proof}
\cref{itm:ch45-inc:psh}.  A continuous function on the compact space
$X^{\An}$ is bounded and upper semicontinuous, and for $j\ge j_0$ we
have $h_{j_0}\le h_j$, so $\mathcal F_\theta(h_j)\supseteq\mathcal F_\theta(h_{j_0})$
is nonempty.  Thus $h_j$ is an admissible obstacle and
\Cref{thm:bj-final-obstacle-psh} applies.  Monotonicity is
\Cref{prop:bj-final-obstacle-calculus}\,(1).

\cref{itm:ch45-inc:sup}.  Each $h_j$ is continuous, hence lower
semicontinuous and bounded below, and $h_j\uparrow h$.
\Cref{lem:obstacle-envelope-formal}\,(4), proved in
\Cref{chap:bj-envelope-history} for raw suprema and without any envelope
input, gives
$\sup_jP^{\mathrm{raw}}_\theta(h_j)=P^{\mathrm{raw}}_\theta(h)$, where
$P^{\mathrm{raw}}_\theta(\cdot)=\sup\mathcal F_\theta(\cdot)$.  By
\Cref{lem:ch45-raw-equals-regularized} the raw and starred envelopes agree
for the admissible obstacles $h_j$, so the left-hand side is
$\sup_jP_\theta(h_j)$.  The right-hand side is $\sup\mathcal F_\theta(h)$ by
definition.

\cref{itm:ch45-inc:star}.  The net $(P_\theta(h_j))_{j\ge j_0}$ is an
upward directed net in $\PSH_X(\theta)$; it is bounded above, because each
of its members lies in $\mathcal F_\theta(h)$, which is bounded above by
hypothesis.  \Cref{thm:bj-final-increasing-nets} therefore applies and gives
a psh regularized supremum, which by \cref{itm:ch45-inc:sup} equals
$\bigl(\sup\mathcal F_\theta(h)\bigr)^*$.  That it is the least upper bound in
$\PSH_X(\theta)$ is \Cref{prop:ch45-lub}\,\cref{itm:ch45-lub:least}.
\end{proof}

\begin{remark}[Why the increasing case is not symmetric to the decreasing one]
\label{rem:ch45-increasing-asymmetry}
\Cref{thm:ch45-decreasing-obstacle-nets} concludes with an equality of
functions; \Cref{thm:ch45-increasing-obstacles} concludes with an equality
only after a star.  The asymmetry is not an artefact.  The limit $h$ of an
increasing net of continuous functions is lower semicontinuous, not upper
semicontinuous, so $h$ is in general \emph{not} an admissible obstacle,
and $\bigl(\sup\mathcal F_\theta(h)\bigr)^*$ need not be $\le h$ --- the
raw supremum always is, every competitor being $\le h$ by
\Cref{def:ch45-admissible-obstacle}\,(3), and it is exactly the
regularization that can escape, as
\Cref{lem:ch45-raw-equals-regularized} and \Cref{rem:ch45-admissibility}
isolate; one cannot even write
$P_\theta(h)$ in the sense of
\Cref{def:ch45-admissible-obstacle}.  What survives is the identity of
suprema in \cref{itm:ch45-inc:sup} and the pshness of the regularization in
\cref{itm:ch45-inc:star}, which is exactly what the applications need.
\end{remark}

\begin{corollary}
\label{cor:ch45-uniform-obstacles}
Let $X$ be integral projective unibranch, $\omega\in\Amp(X)$, and let
$(h_j)\subset C^0(X^{\An})$ be a monotone net --- increasing or
decreasing --- with continuous limit $h\in C^0(X^{\An})$.  Then
$h_j\to h$ uniformly, and
\[
 P_\omega(h_j)\longrightarrow P_\omega(h)
 \qquad\text{uniformly on }X^{\An}.
\]
\end{corollary}

\begin{proof}
The space $X^{\An}$ is compact, so a monotone net of continuous functions
converging pointwise to a continuous limit converges uniformly, by Dini's
theorem for nets recorded in \Cref{app:nets}.  Each $h_j$ and $h$ is an
admissible obstacle for $\omega$, since $\omega$ is nef and the
constants are therefore $\omega$-psh
(\Cref{lem:ch34-psh-nonempty}).  Now
\Cref{prop:ch45-obstacle-calculus-2}\,\cref{itm:ch45-P:lipschitz} gives
\[
 \lVert P_\omega(h_j)-P_\omega(h)\rVert_\infty
 \le\lVert h_j-h\rVert_\infty\longrightarrow0 .
\]
\end{proof}

\subsection{Approximation of psh functions in a fixed class}
\label{sec:ch45-fixed-class-approximation}

\Cref{def:bj-psh} approximates a $\theta$-psh function by generically
finite Fubini--Study functions in classes that merely \emph{converge} to
$\theta$, and \Cref{rem:ch45-real-again} explained why the varying-class
convention is indispensable in general.  For every real ample class,
however, \Cref{thm:ch34-fs-regularization} already gives an unconditional
decreasing net in $\cH_{\mathrm{dom}}(\omega)$, hence a continuous PL
regularization in the fixed class.  The envelope theorem gives a second,
useful construction: the approximants can be indexed by continuous
obstacles and chosen as their psh envelopes.

\begin{theorem}
\label{thm:ch45-continuous-regularization}
Let $X$ be integral projective unibranch, $\omega\in\Amp(X)$, and
$u\in\PSH_X(\omega)$.  Then there is a downward directed net
$(u_g)_{g\in\mathcal G}$ with
\[
 u_g\in C^0(X^{\An})\cap\PSH_X(\omega),
 \qquad
 u_g\downarrow u\ \text{ pointwise on }X^{\An}.
\]
In particular $\PSH_X(\omega)$ is the set of \emph{generically finite}
pointwise limits of decreasing nets of continuous $\omega$-psh functions;
the qualifier is needed, since the constants $-m$ form such a net with
limit $\equiv-\infty$, excluded by \Cref{def:bj-psh}\,(2).
\end{theorem}

\begin{proof}
The function $u$ is upper semicontinuous and bounded above by
$\sup_{X^{\An}}u\in\mathbf R$.  By
\Cref{lem:ch35-usc-inf-continuous} the family
\[
 \mathcal G:=\{g\in C^0(X^{\An})\mid g\ge u\}
\]
is nonempty and downward directed, and $\inf_{g\in\mathcal G}g=u$
pointwise.

For $g\in\mathcal G$ put $u_g:=P_\omega(g)$.  Each $g$ is an
admissible obstacle: it is continuous, hence bounded and upper
semicontinuous, and $u\in\mathcal F_\omega(g)$ by construction.  By
\Cref{cor:ch45-continuity-envelopes}, $u_g\in
C^0(X^{\An})\cap\PSH_X(\omega)$; by
\Cref{prop:bj-final-obstacle-calculus}\,(1) the assignment
$g\mapsto u_g$ is order preserving, so $(u_g)_{g\in\mathcal G}$ is a
downward directed net.

Finally, $u\le u_g\le g$ for every $g$, the first inequality because
$u\in\mathcal F_\omega(g)$ and $u_g$ is the greatest element of that family,
the second by \Cref{thm:bj-final-obstacle-psh}.  Taking the infimum over
$\mathcal G$ gives $u\le\inf_gu_g\le\inf_gg=u$.  Hence $u_g\downarrow
u$.

The last sentence follows: one inclusion is what was just proved, and the
other is \Cref{thm:bj-basic-closure}\,(4), a decreasing net of
$\omega$-psh functions with generically finite limit having
$\omega$-psh limit.
\end{proof}

\begin{remark}[What the fixed class buys, and what it costs]
\label{rem:ch45-fixed-class-cost}
The \emph{existence} asserted here costs no envelope hypothesis in an ample
class: it already follows, for arbitrary projective $X$, from
\Cref{thm:ch34-fs-regularization}.  Ampleness and unibranchness are used by
the present \emph{proof}: ampleness enters through
\Cref{cor:ch45-continuity-envelopes}, while unibranchness supplies the
envelope theorem.  What this construction adds is the distinguished formula
$u_g=P_\omega(g)$ and its compatibility with obstacle calculus.

For a merely nef class the same envelope argument still produces a
decreasing fixed-class net $u_g=P_\theta(g)\downarrow u$, but continuity
of the approximants is no longer guaranteed.  For a pseudoeffective class
that is not nef the construction also remains meaningful, since every
competitor family contains $u$; continuous approximants are then
impossible, because a continuous $\theta$-psh function would be bounded
and would force $\theta\in\Nef(X)$ by
\Cref{lem:ch34-psh-nonempty}.  Outside the ample cone, the varying-class
regularization of \Cref{def:bj-psh}\,(3) is therefore the unconditional
statement, while fixed-class envelope regularization uses the envelope
property.
\end{remark}

\subsection{Conditional joins and meets}
\label{sec:ch45-lattice}

\begin{theorem}
\label{thm:ch45-lattice}
Let $X$ be integral projective unibranch and
$\theta\in N^1(X)_{\mathbf R}$.  Order $\PSH_X(\theta)$ by the pointwise
order on $X^{\An}$.  Let $S\subseteq\PSH_X(\theta)$ be nonempty.
\begin{enumerate}
\item\label{itm:ch45-lat:join}
      If $S$ is bounded above by a real constant, then $S$ has a least
      upper bound in $\PSH_X(\theta)$, namely
      $\bigvee S:=(\sup S)^*$.
\item\label{itm:ch45-lat:meet}
      If $S$ admits at least one lower bound in $\PSH_X(\theta)$, then
      $S$ has a greatest lower bound in $\PSH_X(\theta)$, namely
      $\bigwedge S:=P_\theta(\inf S)$.
\end{enumerate}
In particular, $\max\{u,v\}$ is always the join of $\{u,v\}$.  When
$u,v$ admit a common $\theta$-psh lower bound,
$P_\theta(\min\{u,v\})$ is their meet.  Without that lower-bound
hypothesis, $\PSH_X(\theta)$ need not be a lattice.
\end{theorem}

\begin{proof}
\cref{itm:ch45-lat:join} is \Cref{prop:ch45-lub}.

\cref{itm:ch45-lat:meet}.  Put $f:=\inf S$, an infimum of upper
semicontinuous functions, hence upper semicontinuous: for $c\in\mathbf R$
the set $\{f<c\}=\bigcup_{u\in S}\{u<c\}$ is open.  It is bounded above,
being dominated by any single member of $S$, each of which is bounded
above.  Its competitor family is
\[
 \mathcal F_\theta(f)=\{\psi\in\PSH_X(\theta)\mid\psi\le u\ \text{for all }u\in S\},
\]
the set of $\theta$-psh lower bounds of $S$, which is nonempty by
hypothesis.  So $f$ is an admissible obstacle, and
\Cref{thm:bj-final-obstacle-psh} says that $\mathcal F_\theta(f)$ has a greatest
element, namely $P_\theta(f)$.  A greatest element of the set of lower
bounds is a greatest lower bound.

For the last sentence, $\max\{u,v\}\in\PSH_X(\theta)$ by
\Cref{thm:bj-basic-closure}\,(2) and is evidently the least upper bound;
and $\min\{u,v\}$ is upper semicontinuous and bounded above, so
\cref{itm:ch45-lat:meet} applies to $S=\{u,v\}$ as soon as some
$\theta$-psh function lies below both.
\end{proof}

\begin{remark}[The two conditions are genuinely conditions]
\label{rem:ch45-lattice-conditions}
Neither hypothesis in \Cref{thm:ch45-lattice} can be dropped.

If $S$ is not bounded above by a constant, then $\sup S$ may be
identically $+\infty$, as in \Cref{ex:ch37-pointwise-not-uniform}, and
$S$ has no upper bound in $\PSH_X(\theta)$ at all, a $\theta$-psh
function being bounded above.

If $S$ has no $\theta$-psh lower bound, then it has no greatest one
either, and this occurs.  Take $X=\mathbf P^1$ over an uncountable
algebraically closed field, $\theta=c_1(\cO(1))$, and
$S:=\{\log|s_q|\}_q$ indexed by the closed points $q$, as in
\Cref{rem:ch45-pole-strict}.  A $\theta$-psh lower bound $\psi$ would
satisfy $\psi(v_{q,\triv})\le\log|s_q|(v_{q,\triv})=-\infty$ for every
$q$, that is, $\psi$ would take the value $-\infty$ at uncountably
many ray endpoints.  This is impossible: by the curve criterion recorded at
the end of \Cref{sec:ch34-pl-positivity}, $t\mapsto\psi(t\ord_q)$ is
convex with nondecreasing slope, so $\psi(v_{q,\triv})=-\infty$ forces
$\psi'_q(0^+)<0$; and the criterion also asserts that the sum
$\sum_q\psi'_q(0^+)$, which is bounded below by $-\deg\theta$, has only
countably many nonzero terms.  Hence at most countably many endpoints can
carry the value $-\infty$.
\end{remark}

\section{Geometric characterization in an ample class}
\label{sec:ch45-characterization}

The sufficiency direction is the theorem just proved.
The necessity direction is an earlier Boucksom--Jonsson theorem and is not
part of the supplied envelope proof.
It was stated with provenance in
\Cref{thm:ep-implies-unibranch}.

\begin{theorem}\label{thm:bj-final-unibranch-characterization}
Let $X$ be an integral projective variety over an algebraically closed
trivially valued field.
The following are equivalent:
\begin{enumerate}
\item $X$ is unibranch;
\item $\operatorname{EP}_X(\theta)$ holds for every
      $\theta\in N^1(X)_{\mathbf R}$;
\item $\operatorname{EP}_X(\omega)$ holds for every ample class $\omega$;
\item $\operatorname{EP}_X(\omega)$ holds for at least one ample class
      $\omega$.
\end{enumerate}
\par
\end{theorem}

\begin{proof}
The implication from the first item to the second is
\Cref{thm:bj-final-envelope}.
The second implies the third, and the third implies the fourth.
The fourth implies the first by
\Cref{thm:ep-implies-unibranch}, whose external source is
\cite[Theorem~5.12]{BJ22}.

We expand the four steps, since the cycle is the sharpness statement of
Part~IV and each arrow has a different character.

\emph{(1)$\Rightarrow$(2).}  This is exactly
\Cref{thm:bj-final-envelope}, applied once for each
$\theta\in N^1(X)_{\mathbf R}$; equivalently it is
\Cref{cor:bj-conjecture-514}.  It is the content of nine chapters, and it is
the only implication of the four that is not immediate.

\emph{(2)$\Rightarrow$(3).}  An ample class is an element of
$N^1(X)_{\mathbf R}$, so this is a restriction of the quantifier.

\emph{(3)$\Rightarrow$(4).}  The variety $X$ is projective, so
$\Amp(X)\ne\emptyset$; pick any $\omega$ in it.  This is the only place
where projectivity, as opposed to properness, is used in the cycle.

\emph{(4)$\Rightarrow$(1).}  This is \Cref{thm:ep-implies-unibranch}.
Its mechanism, recalled in \Cref{chap:bj-envelope-history}, is worth keeping
in view because it explains why an ample class is required: the envelope
property for $\omega$ yields, through the birational reflection statement
\cite[Lemma~5.13]{BJ22} applied to the classes $(1+m^{-1})\omega$, that
every function in $\cH_{\mathrm{dom}}(\nu^*\omega)$ descends from
$(X^\nu)^{\An}$ to $X^{\An}$.  Since $\nu$ is finite, $\nu^*\omega$
is ample, so $\cH_{\mathrm{dom}}(\nu^*\omega)$ spans $\PL(X^\nu)$ by the
discussion following \Cref{def:ch34-hdom}; and $\PL(X^\nu)$ separates the
points of $(X^\nu)^{\An}$.  A separating family of functions that all
factor through $\nu^{\An}$ forces $\nu^{\An}$ to be injective, and
non-Archimedean GAGA transfers injectivity to the scheme map.  Ampleness is
used twice: to know that $\nu^*\omega$ is ample, and to know that the
dominating Fubini--Study functions of an ample class span $\PL$.
\end{proof}

\begin{corollary}
\label{cor:ch45-failure-nonunibranch}
Let $X$ be integral projective and not unibranch.  Then:
\begin{enumerate}
\item\label{itm:ch45-fail:ep}
      $\operatorname{EP}_X(\omega)$ fails for \emph{every} ample class
      $\omega$;
\item\label{itm:ch45-fail:compact}
      $\PSH^{\sup}_X(\omega)$ is not compact for the topology of
      pointwise convergence on $X^{\divis}$, for every ample $\omega$;
\item\label{itm:ch45-fail:operator}
      for every ample $\omega$ there is an obstacle
      $h\in C^0(X^{\An})$ for which $P_\omega(h)\notin\PSH_X(\omega)$;
\item\label{itm:ch45-fail:net}
      for every ample $\omega$ there is a bounded-above increasing net in
      $\PSH_X(\omega)$ that does not converge in $\PSH_X(\omega)$.
\end{enumerate}
\end{corollary}

\begin{proof}
\cref{itm:ch45-fail:ep} is the contrapositive of
\Cref{thm:bj-final-unibranch-characterization}\,(4)$\Rightarrow$(1).
\cref{itm:ch45-fail:compact} and \cref{itm:ch45-fail:net} follow from
\cref{itm:ch45-fail:ep} and the equivalences of
\Cref{thm:bj-envelope-equivalences}, which hold for an arbitrary class on an
arbitrary integral projective $X$ and are proved in
\Cref{chap:bj-envelope-history} without any envelope input.
\cref{itm:ch45-fail:operator} follows from \cref{itm:ch45-fail:ep} and
\Cref{thm:general-obstacle-criterion}\,(2), which for an ample class and a
continuous obstacle is again an equivalence; here $\PL(X)\subset
C^0(X^{\An})$ supplies the obstacle.
\end{proof}

\begin{warning}{Necessity is proved only in the ample case}
\label{warn:ch45-necessity-ample}
\Cref{thm:bj-final-unibranch-characterization} characterizes unibranchness
by the envelope property in an ample class.  It does \emph{not} assert that
a variety which is not unibranch fails the envelope property in every class.
Indeed some classes have an empty psh space, for which the envelope property
holds vacuously on any variety whatsoever
(\Cref{rem:ch45-empty-case}); and for a nonzero pseudoeffective class that
is not ample nothing is claimed here in either direction.  The precise
scope is: sufficiency of unibranchness for all real classes
(\Cref{thm:bj-final-envelope}), and necessity of unibranchness for one
ample class (\Cref{thm:ep-implies-unibranch}).  This is also the scope
recorded in \Cref{rem:ch37-not-said}.
\end{warning}

\begin{remark}[The two directions have different provenance]
\label{rem:ch45-two-provenances}
The two implications that make
\Cref{thm:bj-final-unibranch-characterization} a characterization come from
different places and it is a point of principle in this book to say so.
Sufficiency is proved in
\Crefrange{chap:bj-cartier-realization}{chap:bj-envelope-consequences},
following the supplied unpublished manuscript; necessity is
\cite[Theorem~5.12]{BJ22}, published, and is quoted, not proved, here.  The
supplied manuscript states the converse in its introduction as a known
observation of Boucksom and Jonsson, and does not reprove it.  A reader who
wishes to use only published results has
\Cref{thm:ep-implies-unibranch} and the partial results
\Cref{thm:bj22-nef-envelope} and \Cref{thm:bj-addendum-all-classes};
everything beyond those rests on the manuscript.
\end{remark}

\section{Worked examples}
\label{sec:ch45-examples}

Both directions of \Cref{thm:bj-final-unibranch-characterization} are now
available, so the node and the cusp can be computed rather than described.
This section does that, then computes an envelope explicitly on a unibranch
nonnormal curve, and closes with two examples delimiting what the theorem
gives.

\subsection{Curves: the computational conventions}
\label{sec:ch45-curve-conventions}

Throughout this section $C$ is an integral projective curve over the
algebraically closed trivially valued field $k$, with normalization
$\nu:C^\nu\to C$, and $C^\nu$ is a smooth projective curve.  We use the
following three facts, all established earlier.

\emph{(i) The points.}  The divisorial valuations of $k(C)=k(C^\nu)$ are
the $\ord_q$ for $q$ a closed point of $C^\nu$.  For $t\in[0,+\infty)$
the semivaluation $t\ord_q$ is a point of $C^{\An}$ with centre
$\nu(q)$, and the assignment $t\mapsto t\ord_q$ extends to a
homeomorphism of $[0,+\infty]$ onto the closed ray
\[
 R_q:=\{t\ord_q\ :\ t\in[0,+\infty]\}\subset C^{\An},
\]
with $0\ord_q=v_{C,\triv}$ and $+\infty\cdot\ord_q=v_{\nu(q),\triv}$, the
trivial semivaluation of the closed point $\nu(q)$.  Distinct $q$ give
rays meeting only at $v_{C,\triv}$, except that $R_q$ and $R_{q'}$
share the endpoint $v_{\nu(q),\triv}$ when $\nu(q)=\nu(q')$.  This is
the description of \Cref{chap:trivial-analytification} for curves.

\emph{(ii) The criterion on a smooth curve.}  For $C^\nu$ smooth and
$\eta\in N^1(C^\nu)_{\mathbf R}$, a function
$\phi:(C^\nu)^{\An}\to\mathbf R\cup\{-\infty\}$ is $\eta$-psh if and only
if $t\mapsto\phi(t\ord_q)$ is convex on $[0,+\infty)$ for every closed
point $q$ and
\[
 \deg\eta+\sum_q\phi'_q(0^+)\ \ge\ 0 ,
\]
the sum having only countably many nonzero terms.  This is
\cite[Example~4.13]{BJ22}, recorded at the end of
\Cref{sec:ch34-pl-positivity} and checked there against two explicit test
configurations.  The criterion constrains only the restriction of $\phi$
to the open rays; the value at an endpoint is not free, because a
$\theta$-psh function is upper semicontinuous and decreasing, so that
\[
 \phi(v_{\nu(q),\triv})=\lim_{t\to+\infty}\phi(t\ord_q)\in[-\infty,+\infty)
\]
by the argument of \Cref{lem:ch45-branch-symmetry} below.  Accordingly, when
we verify the criterion for an explicitly given $\phi$, we always define
$\phi$ at the endpoints by that limit.  Its existence follows from
monotonicity of a psh function along the ray, not from convexity alone:
a convex function can be increasing.

\emph{(iii) Monotonicity.}  Every $\theta$-psh function on any $X$ is
decreasing for the order on semivaluations: it is a pointwise limit of
generically finite Fubini--Study functions, which are decreasing by
\Cref{prop:ch34-generic-finiteness}\,\cref{itm:ch34-gf:decreasing}, and the
property passes to pointwise limits.  In particular
$u(v_{C,\triv})=\sup_{C^{\An}}u$.

The first lemma is the mechanism behind the node.

\begin{lemma}
\label{lem:ch45-branch-symmetry}
Let $C$ be an integral projective curve, $\theta\in N^1(C)_{\mathbf R}$,
and let $x\in C$ be a closed point.  Let $q,q'$ be closed points of
$C^\nu$ with $\nu(q)=\nu(q')=x$.  Then every $u\in\PSH_C(\theta)$
satisfies
\[
 u(v_{x,\triv})
 =\lim_{t\to+\infty}u(t\ord_{q})
 =\lim_{t\to+\infty}u(t\ord_{q'}) .
\]
In particular the two limits are equal.
\end{lemma}

\begin{proof}
Fix $q$ with $\nu(q)=x$.  For every $t\in[0,+\infty)$ and every
$f\in\cO_{C,x}$ one has $t\ord_q(f)\le v_{x,\triv}(f)$, the right-hand
side being $0$ if $f(x)\ne0$ and $+\infty$ otherwise; hence
$t\ord_q\le v_{x,\triv}$ for the order on semivaluations
\eqref{eq:semivaluation-order}.  By (iii) above, $u$ is decreasing for
that order, so
\[
 u(t\ord_q)\ \ge\ u(v_{x,\triv})
 \qquad\text{for every }t,
\]
and therefore $\lim_{t\to\infty}u(t\ord_q)\ge u(v_{x,\triv})$, the limit
existing in $[-\infty,+\infty)$ because $t\mapsto u(t\ord_q)$ is
nonincreasing.

Conversely, $t\ord_q\to v_{x,\triv}$ in $C^{\An}$ as $t\to+\infty$, by
(i); and $u$ is upper semicontinuous by \Cref{def:bj-psh}\,(1).  Hence
\[
 u(v_{x,\triv})\ \ge\ \limsup_{t\to+\infty}u(t\ord_q)
 =\lim_{t\to+\infty}u(t\ord_q).
\]
The two inequalities give equality.  The same computation with $q'$ in
place of $q$ gives the second equality, both limits being equal to the
single number $u(v_{x,\triv})$.
\end{proof}

\begin{remark}[The lemma is about \texorpdfstring{$C$}{C}, not about \texorpdfstring{$C^\nu$}{C nu}]
\label{rem:ch45-branch-symmetry-scope}
\Cref{lem:ch45-branch-symmetry} is a statement about psh functions on the
possibly singular curve $C$.  Its whole force comes from the fact that the
two rays $R_q$ and $R_{q'}$ of $C^{\An}$ share their endpoint, which
happens exactly when $q$ and $q'$ lie over the same point of $C$.  On
$C^\nu$ the corresponding rays have distinct endpoints
$v_{q,\triv}\ne v_{q',\triv}$, and no such constraint holds; that is the
whole difference between the two spaces, and it is why $\nu^{\An}$ is a
homeomorphism precisely when no two rays are glued.
\end{remark}

\subsection{The node, computed}
\label{sec:ch45-node}

\begin{example}[The nodal cubic]
\label{ex:ch45-node-computed}
Let $\operatorname{char}k\ne2,3$ and let
\[
 C=\{\,y^2z=x^2(x+z)\,\}\subset\mathbf P^2_k
\]
be the integral nodal cubic, with node $x_0=[0:0:1]$.  Its normalization is
\[
 \nu:\ \mathbf P^1\longrightarrow C,
 \qquad
 s\longmapsto[\,s^2-1:\ s(s^2-1):\ 1\,]
\]
on the affine chart, so that $\nu^{-1}(x_0)=\{q_+,q_-\}$ with
$q_\pm=\pm1$, and $\nu$ is an isomorphism elsewhere.  Put
$\omega:=c_1(\cO_C(1))$, an ample class with $\deg\omega=3$; then
$\nu^*\omega=c_1(\cO_{\mathbf P^1}(3))$ has degree $3$ as well.

\emph{The analytic space.}  By \Cref{sec:ch45-curve-conventions}\,(i), the
map $\nu^{\An}$ is a bijection except over $v_{x_0,\triv}$, where it is
two-to-one: the two distinct points $v_{q_+,\triv}$ and $v_{q_-,\triv}$
of $(\mathbf P^1)^{\An}$ both map to $v_{x_0,\triv}$.  Thus $C^{\An}$
is obtained from $(\mathbf P^1)^{\An}$ by gluing the endpoints of the two
rays $R_{q_+}$ and $R_{q_-}$; the union
$R_{q_+}\cup R_{q_-}\subset C^{\An}$ is a loop through $v_{C,\triv}$
and $v_{x_0,\triv}$.  This loop is the analytic trace of the arithmetic
genus of $C$, and it is the geometric reason the two branches cannot be
separated by a function on $C$.

\emph{A psh function upstairs that does not descend.}  Define
$\psi:(\mathbf P^1)^{\An}\to\mathbf R\cup\{-\infty\}$ by
\[
 \psi(t\ord_{q_+})=\max\{-t,-1\}\ \ (t\in[0,+\infty]),
 \qquad
 \psi\equiv0\ \text{off }R_{q_+} .
\]
By the criterion of \Cref{sec:ch45-curve-conventions}\,(ii), $\psi$ is
$\nu^*\omega$-psh: it is convex on every ray, and
$\deg\nu^*\omega+\sum_q\psi'_q(0^+)=3-1=2\ge0$.  Its two branch limits are
\[
 \lim_{t\to\infty}\psi(t\ord_{q_+})=-1,
 \qquad
 \lim_{t\to\infty}\psi(t\ord_{q_-})=0 .
\]
By \Cref{lem:ch45-branch-symmetry}, no $u\in\PSH_C(\omega)$ has
$\nu^*u=\psi$: the pullback of such a $u$ would have equal branch limits
at $q_+$ and $q_-$, both equal to $u(v_{x_0,\triv})$.  Hence
\[
 \nu^*:\ \PSH_C(\omega)\longrightarrow\PSH_{\mathbf P^1}(\nu^*\omega)
\]
is \emph{not} surjective.

\emph{What this shows, and what it does not.}  The failure of surjectivity
is precisely the failure of
\Cref{cor:bj-universal-reflection-all-characteristics} for $\nu$, and it
is exactly what one expects: $\nu$ is finite and birational but not a
universal homeomorphism, since it is not injective.  So the last step of the
assembly in \Cref{sec:ch45-assembly} is unavailable for $C$.  This by
itself does not prove that $\operatorname{EP}_C(\omega)$ fails --- a
different proof of the envelope property might not go through the
normalization --- and that is where the converse is needed:
\Cref{thm:bj-final-unibranch-characterization}\,(4)$\Rightarrow$(1)
gives, since $C$ is not unibranch,
\[
 \operatorname{EP}_C(\omega)\ \text{fails for every ample }\omega .
\]
By \Cref{cor:ch45-failure-nonunibranch} the normalized space
$\PSH^{\sup}_C(\omega)$ is not compact, some continuous obstacle has a
non-psh envelope, and some bounded increasing net of $\omega$-psh
functions fails to converge in $\PSH_C(\omega)$.
\end{example}

\begin{remark}[Where a naive computation on the node goes wrong]
\label{rem:ch45-node-trap}
It is tempting to compute envelopes on $C$ by computing them on
$\mathbf P^1$ and descending, as will be done for the cusp in
\Cref{ex:ch45-cuspidal-envelope}.  \Cref{ex:ch45-node-computed} shows why
this is illegitimate for the node: the two psh classes are not identified,
so a function produced upstairs need not correspond to anything downstairs,
and the largest $\nu^*\omega$-psh function below a pulled-back obstacle
need not be the pullback of the largest $\omega$-psh function below the
obstacle.  Concretely, the obstacle $\nu^*h$ of a continuous $h$ on
$C^{\An}$ automatically satisfies $\nu^*h(v_{q_+,\triv})
=\nu^*h(v_{q_-,\triv})$, but its psh envelope upstairs has no reason to,
and if it does not, it is not a pullback.
\end{remark}

\subsection{The cusp, computed}
\label{sec:ch45-cusp}

\begin{example}[The cuspidal cubic]
\label{ex:ch45-cusp-computed}
Let $\operatorname{char}k\ne2,3$ and let
$C=\{y^2z=x^3\}\subset\mathbf P^2_k$ be the integral cuspidal cubic, with
cusp $x_0=[0:0:1]$.  Its normalization is
\[
 \nu:\ \mathbf P^1\longrightarrow C,
 \qquad
 s\longmapsto[\,s^2:\ s^3:\ 1\,],
\]
a bijection, with $\nu^{-1}(x_0)=\{\tilde p\}$, $\tilde p=0$.  On
completed local rings the inclusion is
\[
 \cO^{\wedge}_{C,x_0}=k[[s^2,s^3]]\ \subset\ k[[s]]
 =\cO^{\wedge}_{\mathbf P^1,\tilde p},
\]
so $C$ is nonnormal and unibranch.  Put
$\omega:=c_1(\cO_C(1))$, $\deg\omega=3$, $\nu^*\omega=c_1(\cO(3))$.

\emph{The identification.}  $\nu$ is finite, birational and bijective, so
by \Cref{lem:bj-unibranch-normalization-universal} it is a universal
homeomorphism and $\nu^{\An}$ is a homeomorphism.  By
\Cref{cor:ch44-unibranch-normalization-reflection}\,\cref{itm:ch44-unib:refl},
\[
 \nu^*:\ \PSH_C(\theta)\xrightarrow{\ \sim\ }\PSH_{\mathbf P^1}(\nu^*\theta)
\]
for every $\theta\in N^1(C)_{\mathbf R}$, preserving the order, finite
maxima, addition of constants and decreasing limits whose limit is not
identically $-\infty$.  By
\Cref{thm:bj-final-envelope} the envelope property holds on $C$ for every
real class.

\emph{Nonnormality is invisible to $\PSH$ but not to sections.}  The
numerical semigroup of the cusp is
\[
 \Gamma:=\{\ord_{\tilde p}(f)\ :\ f\in\cO_{C,x_0}\setminus\{0\}\}
 =\langle2,3\rangle=\{0,2,3,4,5,\dots\},
\]
because $\cO_{C,x_0}$ is generated over $k$ by $u=x/z$ and $w=y/z$
with $\ord_{\tilde p}(u)=2$ and $\ord_{\tilde p}(w)=3$.  The value $1$
is missing.  Consequently there is no section of $\cO_C(m)$ vanishing to
order exactly $1$ at $\tilde p$, and the Fubini--Study function
$\log|s|$ of a single section can never have slope $-1$ along
$R_{\tilde p}$ with $m=1$.  Nevertheless the $\omega$-psh function
with slope $-1$ exists on $C$, by the identification above.  Both
statements are correct, and they illustrate a distinction that the
reflection theorem makes and that is easy to lose: $\nu^*$ is a bijection
of \emph{psh classes}, not of Fubini--Study classes.

\emph{Where the missing slope comes from.}  The resolution is that a
Fubini--Study presentation carries a Veronese denominator, and the local
orders available at the cusp are the elements of $\Gamma$ divided by that
denominator.  For every $m\ge2$ one has $m\in\Gamma$, so
$\cO_{C,x_0}$ contains an element $f$ with
$\ord_{\tilde p}(f)=m$, and a presentation with denominator $m$
reproduces the slope $-1$ at the cusp exactly:
$\tfrac1m\log|f|$ has slope $-\ord_{\tilde p}(f)/m=-1$ along
$R_{\tilde p}$.  So the missing slope is realized exactly, and not merely
approximated, as soon as one is allowed to pass to a multiple; the value
$1$ is missing from $\Gamma$, but $\Gamma$ contains every integer
$\ge2$, so $\Gamma/m$ contains $1$ for every $m\ge2$.  The general
form of this phenomenon is the Frobenius power of
\Cref{lem:bj-uniform-frobenius} in characteristic $p$, and the
additive-power block of \Cref{lem:bj-local-additive-sandwich} in
characteristic zero: in both cases the descent of a Fubini--Study function
along $\nu$ costs a bounded change of the presentation, not a change of
the function.
\end{example}

\begin{remark}[Node versus cusp is a branch count, not a smoothness count]
\label{rem:ch45-branch-not-smooth}
Both curves of \Cref{ex:ch45-node-computed} and
\Cref{ex:ch45-cusp-computed} are singular, both are plane cubics, both have
normalization $\mathbf P^1$, and both are nonnormal at one point.  They
differ in one number, the number of points of the normalization over the
singular point: two for the node, one for the cusp.  That number is exactly
what \Cref{thm:bj-final-unibranch-characterization} detects.  A reader
looking for an invariant that separates the two cases in terms of
singularity type --- multiplicity, delta invariant, embedding dimension ---
will not find one: the node and the cusp have the same delta invariant
$1$ and the same multiplicity $2$.
\end{remark}

\begin{example}[Node and cusp revisited]
\label{ex:ch45-node-cusp-summary}
For an integral nodal curve, the normalization has two branches over the
node.
\Cref{thm:bj-final-unibranch-characterization} shows that the envelope
property fails for every ample class.
For an integral cuspidal curve, the normalization has one branch and is a
universal homeomorphism.
The main theorem gives the envelope property for every real numerical class.
This is a branch distinction, not a smoothness distinction.
\end{example}

\subsection{An envelope computed on a unibranch nonnormal curve}
\label{sec:ch45-cuspidal-envelope}

\begin{example}[An explicit obstacle envelope on the cuspidal cubic]
\label{ex:ch45-cuspidal-envelope}
Keep the data of \Cref{ex:ch45-cusp-computed}: $C$ the cuspidal cubic,
$\nu:\mathbf P^1\to C$ the normalization with
$\nu^{-1}(x_0)=\{\tilde p\}$, and $\omega=c_1(\cO_C(1))$ of degree
$3$.  Define $h:C^{\An}\to\mathbf R$ by
\[
 h(t\ord_{\tilde p}):=-\min\{5t,\,5\}\quad(t\in[0,+\infty]),
 \qquad
 h\equiv0\ \text{ off } R_{\tilde p} .
\]

\emph{$h$ is a continuous obstacle.}  On the ray $R_{\tilde p}$, which
is homeomorphic to $[0,+\infty]$, the formula defines a continuous
function with $h(v_{C,\triv})=0$ and $h(v_{x_0,\triv})=-5$.  Off the ray
$h$ is constant.  Continuity at $v_{C,\triv}$ holds because every
neighbourhood of the trivial valuation contains all but finitely many
complete rays, on which $h=0$, together with an initial segment of
$R_{\tilde p}$, on which $h$ is close to $0$.  Continuity at
$v_{x_0,\triv}$ holds because a neighbourhood of that point meets
$C^{\An}$ only in a terminal segment of $R_{\tilde p}$, on which
$h\equiv-5$.  Since $\omega$ is ample, hence nef, the constants are
$\omega$-psh by \Cref{lem:ch34-psh-nonempty}, and $-5\le h$; so $h$ is
an admissible obstacle in the sense of
\Cref{def:ch45-admissible-obstacle}.

\emph{Transport to $\mathbf P^1$.}  By
\Cref{prop:ch45-obstacle-calculus-2}\,\cref{itm:ch45-P:normalization},
\[
 \nu^*P_\omega(h)=P_{\nu^*\omega}(\nu^*h),
\]
so it suffices to compute the envelope of $\nu^*h$ on $\mathbf P^1$ for
the class $\nu^*\omega$ of degree $3$.  We drop $\nu^*$ from the
notation and write $t$ for the coordinate on $R_{\tilde p}$.

\emph{An upper bound for the competitors.}  Let $u$ be a competitor:
$u\in\PSH_{\mathbf P^1}(\nu^*\omega)$ with $u\le h$.  Put
$a:=u(v_{\triv})=\sup u$, which is finite.  Every slope
$u'_q(0^+)$ is $\le0$, because $u$ is nonincreasing along rays, and
the criterion of \Cref{sec:ch45-curve-conventions}\,(ii) gives
$3+\sum_qu'_q(0^+)\ge0$; hence
\[
 u'_{\tilde p}(0^+)\ \ge\ -3 .
\]
Convexity makes $t\mapsto u(t\ord_{\tilde p})$ have nondecreasing slope,
so $u(t\ord_{\tilde p})\ge a-3t$ for every $t$.  Taking $t=1$ and
using $u\le h$, that is $u(1\cdot\ord_{\tilde p})\le h(1\cdot\ord_{\tilde
p})=-5$, gives
\[
 a-3\ \le\ -5,
 \qquad\text{that is}\qquad
 a\ \le\ -2 .
\]

\emph{The candidate.}  Define $u_0$ by
\[
 u_0(t\ord_{\tilde p}):=\max\{-2-3t,\ -5\}\quad(t\in[0,+\infty]),
 \qquad
 u_0\equiv-2\ \text{ off }R_{\tilde p}.
\]
Then $u_0$ is convex on every ray, being a maximum of two affine functions
on $R_{\tilde p}$ and constant elsewhere; its slopes at the origin are
$-3$ on $R_{\tilde p}$ and $0$ elsewhere, so
\[
 \deg\nu^*\omega+\sum_qu_0'{}_{,q}(0^+)=3-3=0\ \ge\ 0,
\]
and $u_0\in\PSH_{\mathbf P^1}(\nu^*\omega)$ by the criterion.  The case is
extremal: the whole slope budget of the class is spent on one ray.

\emph{$u_0$ is a competitor.}  Off $R_{\tilde p}$ we have
$u_0=-2\le0=h$.  On $R_{\tilde p}$ with $t\in[0,1]$ we have
$h=-5t$, and both $-2-3t\le-5t$ (equivalent to $t\le1$) and
$-5\le-5t$ hold, so $u_0\le h$.  For $t\ge1$ we have $-2-3t\le-5$,
so $u_0=-5=h$.

\emph{$u_0$ dominates every competitor.}  Let $u$ be a competitor, with
$a=u(v_{\triv})\le-2$ as shown.  Off $R_{\tilde p}$, monotonicity gives
$u\le a\le-2=u_0$.  On $R_{\tilde p}$ with $t\in[0,1]$, convexity on
$[0,1]$ gives
\[
 u(t\ord_{\tilde p})\ \le\ (1-t)\,u(0)+t\,u(1\cdot\ord_{\tilde p})
 \ \le\ (1-t)(-2)+t(-5)=-2-3t\ \le\ u_0(t\ord_{\tilde p}).
\]
For $t\ge1$, monotonicity and $u(1\cdot\ord_{\tilde p})\le-5$ give
$u(t\ord_{\tilde p})\le-5=u_0(t\ord_{\tilde p})$.

\emph{Conclusion.}  The raw supremum of the competitor family is $u_0$,
which is already $\nu^*\omega$-psh; hence
\[
 P_{\nu^*\omega}(\nu^*h)=u_0,
 \qquad
 P_\omega(h)=\bigl(\nu^*\bigr)^{-1}(u_0),
\]
the latter being the unique $\omega$-psh function on $C^{\An}$ whose
pullback is $u_0$, which exists by
\Cref{cor:ch44-unibranch-normalization-reflection}\,\cref{itm:ch44-unib:refl}.
Explicitly, on $C^{\An}$,
\[
 P_\omega(h)(t\ord_{\tilde p})=\max\{-2-3t,-5\},
 \qquad
 P_\omega(h)\equiv-2\ \text{ off }R_{\tilde p}.
\]

\emph{Three readings.}  First, the contact set is
\[
 \{P_\omega(h)=h\}=\{t\ord_{\tilde p}\ :\ t\ge1\},
\]
a proper closed subset of $C^{\An}$ containing the proper-support point
$v_{x_0,\triv}$; off it the envelope is strictly below the obstacle, by
$2$ off the ray and by $2-2t$ at $t\in[0,1)$.  Second, the envelope is
\emph{not} obtained by truncating $h$: the value at $v_{C,\triv}$ drops
from $0$ to $-2$, and the drop is forced by the degree of the class
through the slope inequality, not by the obstacle near $v_{C,\triv}$.
Third, the computation used the envelope theorem twice: once through
\Cref{thm:bj-final-obstacle-psh}, to know that a greatest competitor exists
at all, and once through
\Cref{prop:ch45-obstacle-calculus-2}\,\cref{itm:ch45-P:normalization}, to
transport the problem to the normalization.  On a nodal curve neither move
is available, as \Cref{rem:ch45-node-trap} explains.
\end{example}

\begin{remark}[Reading the answer off the degree]
\label{rem:ch45-cuspidal-degree}
The two numbers in \Cref{ex:ch45-cuspidal-envelope} are the two data of the
problem: the slope $-3$ of $P_\omega(h)$ at the origin of the ray is
$-\deg\omega$, and the drop $-2$ at $v_{C,\triv}$ is
$-(5-3)\cdot1$, the excess of the obstacle's slope $5$ over the budget
$3$, multiplied by the length $1$ of the interval on which the obstacle
falls.  Replacing $5$ by any $\lambda>3$ and $1$ by any $T>0$ gives,
by the same computation,
\[
 P_\omega(h_{\lambda,T})(t\ord_{\tilde p})
 =\max\{-(\lambda-3)T-3t,\ -\lambda T\},
\]
with contact set $\{t\ge T\}$; and for $\lambda\le3$ the obstacle is
already $\omega$-psh and $P_\omega(h_{\lambda,T})=h_{\lambda,T}$, by
\Cref{prop:ch45-obstacle-calculus-2}\,\cref{itm:ch45-P:projection}.  The
transition at $\lambda=3$ is exactly the transition between an obstacle
that the class can support and one that it cannot.
\end{remark}

\subsection{What the theorem does not give}
\label{sec:ch45-not-given}

\begin{example}[No attainment]
\label{ex:ch45-no-attainment}
Let $X=\mathbf P^1$, $\theta=c_1(\cO(1))$ of degree $1$, and let
$p\in X$ be a closed point.  For $c\in(0,1)$ define
\[
 \varphi_c(t\ord_p):=\max\{-t,\,-c\}\ \ (t\in[0,+\infty]),
 \qquad
 \varphi_c\equiv0\ \text{ off }R_p .
\]
Each $\varphi_c$ is $\theta$-psh by the criterion of
\Cref{sec:ch45-curve-conventions}\,(ii): it is convex on every ray and
$1+(-1)=0\ge0$.  The family is bounded above by $0$ and is increasing as
$c$ decreases, with
\[
 \sup_{c\in(0,1)}\varphi_c\ \equiv\ 0
\]
pointwise on all of $X^{\An}$ --- including at $v_{p,\triv}$, where
$\varphi_c(v_{p,\triv})=-c\to0$.  The supremum is therefore already upper
semicontinuous, and \Cref{thm:bj-final-envelope} returns the constant $0$,
which is indeed $\theta$-psh.

But $0\ne\varphi_c$ for every $c$, since $\varphi_c(v_{p,\triv})=-c<0$.
So the least upper bound of the family in $\PSH_X(\theta)$ is not a member
of the family, and the same holds for every subfamily \emph{with the same
supremum}: such a subfamily has parameters accumulating at $0$, hence no
greatest element.  (A singleton subfamily of course has one, so the
qualification cannot be dropped.)  The envelope theorem
produces a supremum, never a maximum.
\end{example}

\begin{example}[No regularity: a continuous, non-piecewise-linear envelope]
\label{ex:ch45-no-regularity}
Let $X=\mathbf P^1$, $\theta=c_1(\cO(1))$ of degree $1$, and
$p\in X$ a closed point.  Define $h\in C^0(X^{\An})$ by
\[
 h(t\ord_p):=-2\bigl(1-\rme^{-t}\bigr)\ \ (t\in[0,+\infty]),
 \qquad
 h\equiv0\ \text{ off } R_p ,
\]
so that $h(v_{p,\triv})=-2$.  Continuity is checked as in
\Cref{ex:ch45-cuspidal-envelope}.  The obstacle is admissible: $\theta$ is
ample, so the constant $-2\le h$ is a competitor.

\emph{The envelope.}  Set $\tau:=\log2$.  We claim
\[
 P_\theta(h)(t\ord_p)=
 \begin{cases}
  \tau-1-t, & 0\le t\le\tau,\\[2pt]
  -2\bigl(1-\rme^{-t}\bigr), & \tau\le t\le+\infty,
 \end{cases}
 \qquad
 P_\theta(h)\equiv\tau-1\ \text{ off } R_p .
\]
Write $u_0$ for the right-hand side.  The two formulas agree at
$t=\tau$, both giving $-1$, and their one-sided derivatives there are
$-1$ and $-2\rme^{-\tau}=-1$; so $u_0$ is of class $C^1$ along
$R_p$, with slope $-1$ on $[0,\tau]$ and slope
$-2\rme^{-t}\in[-1,0)$ increasing on $[\tau,+\infty)$.  Hence $u_0$ is
convex on every ray, and
$\deg\theta+\sum_qu_0'{}_{,q}(0^+)=1-1=0\ge0$; so
$u_0\in\PSH_X(\theta)$.

\emph{$u_0$ is a competitor.}  Off $R_p$, $\tau-1<0=h$.  On $R_p$
with $t\ge\tau$ one has $u_0=h$.  For $t\in[0,\tau]$ the inequality
$u_0\le h$ reads $\tau-1\le g(t):=t-2+2\rme^{-t}$; and
$g'(t)=1-2\rme^{-t}$ vanishes only at $t=\tau$, where $g$ attains its
minimum $g(\tau)=\tau-1$.  So $u_0\le h$ with equality on
$\{t\ge\tau\}$.

\emph{Maximality.}  Let $u$ be a competitor and $a:=u(v_{\triv})$.  As
in \Cref{ex:ch45-cuspidal-envelope}, $u'_p(0^+)\ge-1$ and convexity give
$u(t\ord_p)\ge a-t$; combined with $u\le h$ this forces
$a\le g(t)$ for every $t\ge0$, hence $a\le\tau-1$.  Off $R_p$,
monotonicity gives $u\le a\le\tau-1=u_0$.  On $R_p$ with
$t\in[0,\tau]$, convexity on $[0,\tau]$ gives
\[
 u(t\ord_p)\le\Bigl(1-\tfrac t\tau\Bigr)a+\tfrac t\tau\,u(\tau\ord_p)
 \le\Bigl(1-\tfrac t\tau\Bigr)(\tau-1)+\tfrac t\tau(-1)
 =\tau-1-t=u_0(t\ord_p),
\]
using $u(\tau\ord_p)\le h(\tau\ord_p)=-1$; and for $t\ge\tau$,
$u\le h=u_0$.  Hence $P_\theta(h)=u_0$.

\emph{What the example shows.}  The class is ample and the obstacle is
real analytic along the ray, yet the envelope is only $C^1$: its
second derivative jumps at $t=\tau$, from $0$ to $2\rme^{-\tau}=1$.
It is not a PL function, not a Fubini--Study function, and the constant
$\tau-1=\log2-1$ is transcendental.  So no algebraic or piecewise-linear
regularity can be expected of $P_\theta$, and
\Cref{cor:ch45-continuity-envelopes} is sharp: continuity is what the
ample case guarantees, without a general piecewise-linear or $C^2$
conclusion.
\end{example}

\begin{remark}[Three further things the theorem does not give]
\label{rem:ch45-not-given-list}
Beyond \Cref{ex:ch45-no-attainment} and \Cref{ex:ch45-no-regularity}:
\begin{enumerate}
\item \emph{No finiteness.}  The envelope may be $-\infty$ somewhere, by
      \Cref{prop:ch45-envelope-can-be-minus-infinity}; and by
      \Cref{cor:ch45-extremal} the envelope of the continuous obstacle $0$
      is unbounded below whenever $\PSH_X(\theta)\ne\emptyset$ and
      $\theta$ is not nef.
\item \emph{No continuity outside the ample case.}
      \Cref{cor:ch45-continuity-envelopes} is stated for ample classes and
      \Cref{warn:ch45-continuity-ample} explains why; outside the ample case,
      when the competitor family is nonempty, the psh envelope of a
      continuous obstacle need not be continuous.
\item \emph{No quantitative comparison of $\sup_i\varphi_i$ with its
      regularization.}  The theorem asserts that the regularization is psh;
      it says nothing about the size of the set where the two differ, beyond
      the fact --- \Cref{thm:bj-final-increasing-nets} --- that they agree on
      $X^{\divis}$.  In \Cref{ex:ch45-raw-sup-not-usc} they differ at one
      point; there is no statement here bounding that locus in general.
\end{enumerate}
\end{remark}

\section{Failure boundaries}
\label{sec:ch45-failure}

\Cref{prop:bj-final-hypothesis-ledger} recorded the exact scope of the
hypotheses.  Here we record what happens outside them, so that the theorem
is not quoted where it does not apply.

\subsection{Dropping unibranchness}
\label{sec:ch45-failure-unibranch}

This is the only hypothesis for which a converse is available, and
\Cref{cor:ch45-failure-nonunibranch} states the failure precisely: on an
integral projective $X$ that is not unibranch, the envelope property fails
in \emph{every} ample class, together with its compactness, obstacle and
increasing-net reformulations.  Nothing is claimed for non-ample classes on
such an $X$; see \Cref{warn:ch45-necessity-ample}.

\subsection{Dropping integrality}
\label{sec:ch45-failure-integral}

The class $\PSH_X(\theta)$ is defined in \Cref{chap:bj-fs-pl-psh} for a
projective $X$ with the generic finiteness condition
\Cref{def:bj-psh}\,(2) imposed componentwise, so the envelope property makes
sense for a reducible $X$.  Nothing in this book proves it there.  Two
remarks delimit the gap.  First, on a disconnected $X$ the property is
componentwise, since both $\PSH$ and the regularized supremum are formed
componentwise; so the substantive case is $X$ connected and reducible.
Second, on a connected reducible $X$ the mechanism of
\Cref{thm:ep-implies-unibranch} still shows that the envelope property in an
ample class forces $\nu^{\An}$ to be injective, hence forces $X$ to be
irreducible; so on a connected reducible $X$ the envelope property already
fails in the ample case.  What is not addressed is the behaviour of the
non-ample classes, and the whole of the sufficiency direction, whose proof
uses integrality at every step --- b-divisors, alterations, normalization
--- and has not been extended to reducible schemes here.

\subsection{Dropping projectivity}
\label{sec:ch45-failure-projective}

Projectivity is used in four places and each use is essential to the proof
as written.  It supplies compactness of $X^{\An}$, without which
``bounded above'' loses its force; it supplies the ample class in which the
alterations of \Cref{chap:bj-alterations} and the Stein factorization
operate; it supplies the semiample approximations in the Cartier realization
of \Cref{chap:bj-cartier-realization}; and it supplies the very definition
of $N^1(X)_{\mathbf R}$ and of Fubini--Study functions, which are built
from global sections.  For a proper non-projective $X$ the statement can
be formulated by Chow's lemma but is not proved here.

\subsection{Dropping the trivial absolute value}
\label{sec:ch45-failure-trivial}

The whole of Part~IV is set over a trivially valued $k$, and the objects
are not the same over a nontrivially valued field: there the relevant theory
is the BFJ theory of \Cref{part:bfj}, whose envelope statements
(\Cref{chap:bfj-envelopes}) are proved by different means and under
different hypotheses.  No statement of this chapter is asserted over a
nontrivially valued base, and no statement of \Cref{chap:bfj-envelopes} is
used here.

\subsection{Dropping algebraic closedness}
\label{sec:ch45-failure-closed}

Algebraic closedness of $k$ is used in two distinct ways.  It makes
regularity equivalent to smoothness, which is what allows the de Jong
alteration of \Cref{thm:bj-de-jong-input} to feed
\Cref{thm:bj-smooth-envelope}; and it is what makes a finite birational
bijective morphism radicial in the proof of
\Cref{lem:bj-unibranch-normalization-universal}, through the diagonal
criterion and the $k$-rationality of closed points.  Over a
non-algebraically-closed $k$ both steps require modification, and neither
modification is carried out here.

\subsection{What is not claimed}
\label{sec:ch45-failure-notclaimed}

No claim is made that an arbitrary raw supremum is psh without upper
semicontinuous regularization; \Cref{ex:ch45-raw-sup-not-usc} refutes it.
No claim is made that the supremum is attained;
\Cref{ex:ch45-no-attainment} refutes it.  No claim is made about
regularity; \Cref{ex:ch45-no-regularity} shows that even in the best case
--- ample class, real analytic obstacle --- the envelope is only $C^1$ and
involves transcendental constants.  No claim is made that
$P_\theta$ is continuous for the topology of pointwise convergence, nor
that it commutes with arbitrary limits: what is proved is
\Cref{thm:ch45-decreasing-obstacle-nets} for decreasing nets and
\Cref{thm:ch45-increasing-obstacles} for increasing nets of continuous
obstacles, and those are the two monotone cases.

\section{Chapter dictionary}
\label{sec:ch45-dictionary}

\begingroup
\begin{longtable}{@{}p{.23\textwidth}p{.70\textwidth}@{}}
\toprule
statement & conclusion and scope\\
\midrule
\endfirsthead
\toprule
statement & conclusion and scope\\
\midrule
\endhead
the assembly &
Regularized suprema commute with normalization pullback; the envelope
property upstairs and psh reflection give it downstairs.  Here $X$ is
integral projective unibranch, $k$ is algebraically closed and trivially
valued, and $\theta$ is real
(\Cref{thm:bj-final-envelope,sec:ch45-assembly}).\\
non-circularity &
The six assembly-input groups are established earlier in Part~IV; the
published converse is used only for the characterization
(\Cref{prop:ch45-noncircular}).\\
increasing nets &
A bounded-above increasing net of psh functions has a psh regularized
supremum with the expected divisorial trace.  No sequential reduction is
assumed (\Cref{thm:bj-final-increasing-nets,warn:ch45-nets}).\\
least upper bound &
Regularization computes the join in $\PSH_X(\theta)$ for each nonempty
family bounded above by a constant (\Cref{prop:ch45-lub}).\\
compactness &
The sup-normalized psh space is compact in the divisorial pointwise
topology exactly when the envelope property holds; it is compact for
unibranch $X$ (\Cref{cor:ch45-compactness,rem:ch45-compactness-equivalent}).\\
perturbation of the class &
The envelope property in big rational classes implies it in all real
classes by approximation from above with nef class gaps
(\Cref{prop:ch45-reduction-big-rational,rem:ch45-two-orientations}).\\
obstacle envelope &
For an upper semicontinuous, bounded-above obstacle with a psh competitor,
$P_\theta(f)$ is the greatest such competitor; the raw and regularized
definitions agree
(\Cref{def:ch45-admissible-obstacle,thm:bj-final-obstacle-psh,lem:ch45-raw-equals-regularized}).\\
calculus of $P_\theta$ &
On admissible obstacles the operator is monotone, constant-equivariant,
idempotent, determined by divisorial data, and compatible with
normalization; on bounded real obstacles it is $1$-Lipschitz
when $\theta$ is nef
(\Cref{prop:bj-final-obstacle-calculus,prop:ch45-obstacle-calculus-2}).\\
extremal function &
When $\PSH_X(\theta)$ is nonempty, $V_\theta=P_\theta(0)$ is zero
exactly for nef $\theta$ (\Cref{cor:ch45-extremal}).\\
monotone obstacles &
Decreasing nets commute with $P_\theta$ under a common competitor;
increasing continuous obstacles commute after regularization under the
stated boundedness hypotheses.  In an ample class, continuous monotone
limits give uniform convergence of envelopes
(\Cref{thm:ch45-decreasing-obstacle-nets,thm:ch45-increasing-obstacles,cor:ch45-uniform-obstacles}).\\
conditional joins and meets &
On unibranch $X$, joins exist for constant-bounded families; meets exist
when the family has a psh lower bound (\Cref{thm:ch45-lattice}).\\
geometric characterization &
For integral projective $X$, unibranchness is equivalent to the envelope
property in one ample class.  Necessity is the separate published result
\cite[Theorem~5.12]{BJ22}
(\Cref{thm:bj-final-unibranch-characterization}).\\
node versus cusp &
On an integral projective curve, psh branch limits agree at a shared
endpoint.  A node has two such branches; a cusp has one
(\Cref{lem:ch45-branch-symmetry,ex:ch45-node-computed,ex:ch45-cusp-computed}).\\
\bottomrule
\end{longtable}
\endgroup

Four separations should be carried forward.

\emph{Property versus function.}  The envelope property is a statement about
all bounded-above families at once; the reflection theorem is a statement
about individual functions.  The assembly needs both, and they enter at
different steps (\Cref{sec:ch45-attributions},
\Cref{sec:ch45-four-steps}).

\emph{Supremum versus maximum.}  \Cref{thm:bj-final-envelope} produces a
least upper bound, which is a member of $\PSH_X(\theta)$ but need not be a
member of the given family (\Cref{ex:ch45-no-attainment}).  For obstacle
envelopes the situation is better: $P_\theta(f)$ \emph{is} a member of
$\mathcal F_\theta(f)$, and that is exactly what upper semicontinuity of the
obstacle buys (\Cref{lem:ch45-raw-equals-regularized}).

\emph{Raw versus regularized.}  Before this chapter the two obstacle
envelopes of \Cref{def:bj-obstacle-envelope} and
\Cref{def:bj-final-obstacle-envelope} were different objects and the
notation $P_\theta$ was ambiguous.  They agree for admissible obstacles;
without admissibility this is no longer guaranteed.
\Cref{rem:ch45-admissibility} exhibits an obstacle for which the starred
version is not even a competitor.

\emph{Closure versus regularity.}  The envelope theorem itself is a
closure statement.  Continuity, in the ample-class case, comes
from the separate criterion \eqref{eq:continuity-of-envelopes} for ample
classes, and even then no better regularity is guaranteed
(\Cref{ex:ch45-no-regularity}).

\section{Historical replacement and provenance}
\label{sec:ch45-historical}

The first supplied book treated the envelope property as
Conjecture~5.14 of \cite{BJ22}.
For integral projective unibranch varieties, that paragraph is now replaced
by \Cref{thm:bj-final-envelope} and its proof in
\Crefrange{chap:bj-envelope-statement}{chap:bj-envelope-consequences}.
The theorem covers arbitrary dimension, arbitrary characteristic, every real
numerical class, and all proper-support values.

The proof of the sufficiency theorem and the assembly
\eqref{eq:bj-final-pull-envelope} come from the supplied manuscript
\emph{The envelope property for unibranch projective varieties}.
The increasing-net formulation is \cite[Theorem~5.11]{BJ22}.
The necessity of unibranchness in an ample class is the separate result
\cite[Theorem~5.12]{BJ22}, already recorded in
\Cref{chap:bj-envelope-history}.
This provenance split is important:
the supplied envelope source proves the unibranch sufficiency direction but
does not itself supply the converse proof.

No claim is made here for reducible schemes, nonprojective spaces, or
nontrivially valued ground fields.
No claim is made that an arbitrary raw supremum is psh without upper
semicontinuous regularization.
Those are genuine boundaries of the theorem as proved.

\subsection{The conjecture, and what replaces it}
\label{sec:ch45-conjecture-replaced}

The statement being replaced is the following, quoted in the form in which
it appears in the source.

\begin{quote}
\emph{Conjecture 5.14 of \cite{BJ22}.}  If $X$ is unibranch, then the
envelope property holds for all $\theta\in N^1(X)$.
\end{quote}

Three features of the conjecture should be noted before it is retired.
First, it is stated for an integral projective variety over an algebraically
closed trivially valued field, which is the standing setting of that paper
and of Part~IV; no extension to reducible or nonprojective $X$ is
conjectured there.  Second, it asserts the property for \emph{all} classes,
not merely ample ones; this is why \Cref{thm:bj-final-envelope} is stated
for an arbitrary $\theta\in N^1(X)_{\mathbf R}$ and why
\Cref{sec:ch45-perturbation} takes the reduction of classes seriously.
Third, its converse in the ample case was already a theorem when the
conjecture was formulated, namely \cite[Theorem~5.12]{BJ22}; so the
conjecture and that theorem were designed to fit together into the
characterization that \Cref{thm:bj-final-unibranch-characterization} now
records.

What replaces it in this book is:

\begin{itemize}
\item \Cref{thm:bj-final-envelope} and \Cref{cor:bj-conjecture-514}, the
      conjecture itself, proved;
\item \Cref{thm:bj-final-unibranch-characterization}, the resulting
      characterization of unibranchness by the envelope property in an
      ample class;
\item \Cref{cor:ch45-failure-nonunibranch}, the failure statement outside
      the unibranch case, which is the practical form of the converse.
\end{itemize}

Every occurrence of the phrase ``conjecturally'' attached to the envelope
property in the earlier literature, and every hypothesis of the form
``assume the envelope property'' in
\Crefrange{chap:bj-energy}{chap:bj-stability}, is discharged for unibranch
$X$ by these statements.  The hypothesis is retained in the statements of
those chapters where the underlying source states it, because the source's
theorem is the conditional one and the book's policy is to state the theorem
that was proved and then discharge its hypothesis; it is not retained
because there is any doubt about the discharge.

\subsection{The line of results}
\label{sec:ch45-line-of-results}

The following is the chain of statements this chapter completes.

\begin{enumerate}
\item \emph{The complex prototype.}  Upper envelopes of plurisubharmonic
      functions on a domain in $\mathbf C^n$ are plurisubharmonic after
      upper semicontinuous regularization; this is classical, and the
      version with the regularization made explicit is Bedford--Taylor
      \cite{BT82}.  In the complex analytic setting the corresponding
      statement on a possibly singular space holds when the space is locally
      irreducible, that is, unibranch; \cite[\S5.2]{BJ22} records this as
      the motivation for Conjecture~5.14.
\item \emph{Nef classes on smooth varieties.}  \cite[Theorem~5.20]{BJ22}
      proves the envelope property for every nef class on a smooth connected
      $X$, under the alternative hypothesis
      $\operatorname{char}k=0$ or $\dim X\le2$.  This is
      \Cref{thm:bj22-nef-envelope}.
\item \emph{All classes on smooth varieties.}  The addendum
      \cite{BJ22b} removes the nefness restriction under the same
      alternative hypothesis; this is \Cref{thm:bj-addendum-all-classes},
      Theorem~A of that paper.
\item \emph{All classes on all unibranch varieties, all characteristics.}
      This is \Cref{thm:bj-final-envelope}, from the supplied manuscript.
      The three restrictions removed relative to the previous item are:
      smoothness of $X$, weakened to unibranchness; the
      characteristic-or-dimension alternative, removed entirely; and the
      implicit restriction to the classes for which the earlier obstacle
      machinery applied.
\item \emph{The converse.}  \cite[Theorem~5.12]{BJ22}, that the envelope
      property in one ample class forces $X$ to be unibranch; this is
      \Cref{thm:ep-implies-unibranch} and is not part of the supplied
      manuscript.
\end{enumerate}

The single technical innovation that makes step~(4) possible in positive
characteristic is the replacement of asymptotic multiplier ideals by
asymptotic test ideals, together with Canton's valuative bound
\cite{Canton} and a relative form of Musta\c t\u a's uniform generation
theorem \cite{MustataTest}; this is
\Cref{chap:bj-smooth-characteristic}.  The single innovation that makes the
passage from normal to unibranch possible is the treatment of the
normalization as a universal homeomorphism rather than as a finite or
birational map: Frobenius in characteristic $p$
(\Cref{chap:bj-universal-homeomorphisms}) and the Roberts--Singh
additive-power theorem \cite{RobertsSingh} in characteristic zero
(\Cref{chap:bj-subintegral-normalization}).

\subsection{What becomes unconditional}
\label{sec:ch45-unconditional}

A number of statements in the literature on trivially valued pluripotential
theory carry the envelope property as an explicit hypothesis.  On an
integral projective unibranch $X$ they are now unconditional.  Within this
book the relevant ones are the following, and each is quoted in its own
chapter with the hypothesis displayed and then discharged by
\Cref{thm:bj-final-envelope}:

\begin{itemize}
\item the equivalences of \Cref{thm:bj-envelope-equivalences} and
      \Cref{thm:general-obstacle-criterion}, and the continuity criterion
      \eqref{eq:continuity-of-envelopes} for ample classes, all in
      \Cref{chap:bj-envelope-history};
\item the differentiability and orthogonality statements of
      \Cref{chap:bj-energy} that are proved by envelope perturbation;
\item the existence of maximizers, and hence the non-Archimedean
      Calabi--Yau theorem \Cref{cor:bj-cy-unibranch}, in
      \Cref{chap:bj-finite-energy};
\item the surjectivity of the Monge--Amp\`ere operator onto finite-energy
      measures and the identification of negligible with pluripolar sets in
      \Cref{chap:bj-strong-topologies};
\item the completeness of $\cE^1(\omega)$ and the valuative criteria in
      \Cref{chap:bj-stability}.
\end{itemize}

We do not assert a blanket statement of the form ``every conditional result
in the literature is now unconditional''.  Such a statement cannot be
checked and would be worthless if it could not; each application is
discharged individually, at the place where its hypothesis is stated.

\section{Source and provenance}
\label{sec:ch45-source-note}

\subsection{Correspondence with the supplied manuscript}
\label{sec:ch45-manuscript-correspondence}

This chapter corresponds to \S7 of the supplied manuscript \emph{The
envelope property for unibranch projective varieties}, together with the
parts of its \S2 that concern the reformulations of the envelope property.
The correspondence is:

\begin{center}
\begin{tabular}{p{.34\textwidth}p{.56\textwidth}}
\toprule
manuscript & this chapter\\
\midrule
Theorem 7.1 (= Theorem A) &
\Cref{thm:bj-final-envelope}, restated in \Cref{chap:bj-envelope-statement}
and proved in \Cref{sec:ch45-assembly}\\
the displayed chain smooth $\Rightarrow$ normal $\Rightarrow$ unibranch
of its \S1.2 &
\eqref{eq:bj-final-dependency-diagram} and the attributed chain of
\Cref{sec:ch45-attributions}\\
Lemma 6.8 &
\Cref{lem:bj-unibranch-normalization-universal}, proved in
\Cref{chap:bj-universal-homeomorphisms}; recalled here as
\Cref{lem:ch45-normalization-data}\,\cref{itm:ch45-nd:uh}\\
Theorem 6.7 &
\Cref{cor:bj-universal-reflection-all-characteristics}, proved in
\Crefrange{chap:bj-universal-homeomorphisms}{chap:bj-subintegral-normalization};
recalled here as
\Cref{lem:ch45-normalization-data}\,\cref{itm:ch45-nd:refl}\\
Corollary 5.6 &
\Cref{thm:bj-normal-envelope} and
\Cref{cor:bj-normalization-source-ep}, proved in
\Cref{chap:bj-alterations}\\
its equation (7.1) &
\eqref{eq:bj-final-pull-envelope}\\
its Remark 7.2 (exactness at proper-support points) &
\Cref{sec:ch45-endpoint}, \Cref{prop:bj-final-no-hidden-finiteness}\\
its Remark 7.3 (conditional statements become unconditional) &
\Cref{sec:ch45-unconditional}\\
Proposition 2.3 (increasing-net form) &
\Cref{prop:bj-ep-equivalences} and
\Cref{thm:bj-final-increasing-nets}\\
the perturbation lemma quoted in its \S2.2 and \S2.5 &
\Cref{lem:ch40-ep-nef-limit},
\Cref{lem:ch40-rational-approximation},
\Cref{prop:ch45-reduction-big-rational}\\
\bottomrule
\end{tabular}
\end{center}

The manuscript is an unpublished source supplied for this project.  Results
proved only there are identified as such throughout Part~IV and are never
presented as previously published theorems; the policy is stated in the
front matter and applied here without exception.
\Cref{app:provenance} records the version and date.

\subsection{What is supplied here and is not in the source}
\label{sec:ch45-additions}

The following material of this chapter is not in the manuscript in this
form and is supplied here.

\Cref{prop:ch45-noncircular} and \Cref{rem:ch45-audit-scope}: the
acyclicity audit.  The manuscript's proof is linear by construction and does
not need an audit; a book that spends nine chapters on the proof does.

\Cref{lem:ch45-normalization-data} and \Cref{sec:ch45-four-steps}: the
hypothesis-by-hypothesis verification of the assembly.

\Cref{prop:ch45-lub}, \Cref{cor:ch45-compactness},
\Cref{prop:ch45-reduction-big-rational},
\Cref{cor:ch45-big-rational-suffices}: the reformulations, assembled from
statements that exist separately in \cite{BJ22} and in
\Cref{chap:bj-smooth-characteristic}.

\Cref{def:ch45-admissible-obstacle},
\Cref{lem:ch45-raw-equals-regularized},
\Cref{prop:ch45-obstacle-calculus-2},
\Cref{cor:ch45-extremal},
\Cref{prop:ch45-trichotomy},
\Cref{cor:ch45-continuity-envelopes}: the obstacle calculus.  The
manuscript uses the obstacle envelope only on smooth varieties and only as
an intermediate device, so its calculus is not developed there; the raw
formal identities are \Cref{lem:obstacle-envelope-formal}, quoted in
\Cref{chap:bj-envelope-history} from \cite[\S5.3]{BJ22}, and what is new is
that they now hold for an operator with values in $\PSH_X(\theta)$.

\Cref{thm:ch45-decreasing-obstacle-nets},
\Cref{thm:ch45-increasing-obstacles},
\Cref{cor:ch45-uniform-obstacles},
\Cref{thm:ch45-continuous-regularization},
\Cref{thm:ch45-lattice}: the monotone calculus and the conditional join--meet
structure.

\Cref{lem:ch45-branch-symmetry} and the computations
\Cref{ex:ch45-node-computed},
\Cref{ex:ch45-cusp-computed},
\Cref{ex:ch45-cuspidal-envelope},
\Cref{ex:ch45-raw-sup-not-usc},
\Cref{ex:ch45-no-attainment},
\Cref{ex:ch45-no-regularity}: all new.  They are included because the
chapter's two claims --- that the theorem is exactly as strong as
unibranchness, and that it is a closure statement and nothing more --- are
exactly the claims that a plausible reading can appear to establish without
doing so.  A reader who has computed the branch limits on the node, the
numerical semigroup on the cusp, and the constant $\log2-1$ in
\Cref{ex:ch45-no-regularity} will not confuse a closure theorem with a
regularity theorem.

\subsection{External attributions}
\label{sec:ch45-external}

The envelope property was formulated, and its two reformulations proved, in
\cite[\S5]{BJ22}: the increasing-net and compactness equivalences are
Theorem~5.11 there, the comparison interface used throughout is
Theorem~4.22, the perturbation lemma is Lemma~5.9, the obstacle envelope and
its formal identities are \S5.3, the birational reflection lemma used in the
converse is Lemma~5.13, and the converse itself is Theorem~5.12.  The
positive partial results are Theorem~5.20 there and Theorem~A of the
addendum \cite{BJ22b}.  The complex prototype is \cite{BT82}.  The curve
criterion used in \Cref{sec:ch45-curve-conventions} is Example~4.13 of
\cite{BJ22}.  Everything else in this chapter is either proved here or
proved in
\Crefrange{chap:bj-cartier-realization}{chap:bj-subintegral-normalization}
from the supplied manuscript.

\subsection{Historical note}
\label{sec:ch45-history-note}

The envelope property entered the trivially valued theory as the
non-Archimedean shadow of a classical fact about plurisubharmonic functions,
and it was expected to be soft.  It is not.  The reason, visible in
retrospect in \Cref{lem:ch45-branch-symmetry}, is that the analytification
of a singular variety over a trivially valued field remembers the branch
structure of the singularity in a way that the Zariski topology does not: two
branches through a point produce two rays glued at a single endpoint, and a
function on the analytification cannot separate them, while a function on
the normalization can.  The envelope property is precisely the statement
that this is the only obstruction.

The proof took the form it has because the three natural classes of
morphisms behave differently.  Birational morphisms are handled by
b-divisors, which is the classical route and requires the realization
theorem of \Cref{chap:bj-cartier-realization} to avoid presupposing what is
to be proved.  Finite morphisms are handled by averaging, which requires a
norm that is a universal polynomial law rather than a product over sheets,
so that inseparability and nonreduced fibres cause no loss
(\Cref{chap:bj-finite-descent}).  The normalization of a unibranch variety
is both, and neither: it is the degenerate case in which the previous two
techniques have nothing to bite on, and it is handled by the observation
that it is a universal homeomorphism, so that the only thing to descend is
the algebra.  That last observation is what makes the characteristic-free
statement possible, and it is why the two characteristics receive genuinely
different arguments --- Frobenius and Roberts--Singh --- rather than a single
argument with a case distinction.

Finally, a word on the status of the result in this book.  The sufficiency
direction rests on an unpublished manuscript supplied for this project.  It
is reproduced here in full, with every hypothesis checked and every external
input stated in the exact form used, so that a reader can verify it without
the manuscript; but its provenance is stated wherever it is used, and no
statement of this book converts it into a published theorem by citation.
The necessity direction is published and is quoted as such.  The
characterization \Cref{thm:bj-final-unibranch-characterization} is therefore
a hybrid, and the two halves are labelled.
